\section{Critical levels, exact Hausdorff gauges, the positive-level trichotomy, and general recurrences}\label{sec:gauge}

\subsection{Critical levels and exact gauges}

For an integer $k\ge1$, set $b=2k+1$. Define an increasing continuous
gauge $h_k$ by $h_k(0)=0$, $h_k(1/Q_n)=b^{-n}$ ($n\ge1$), and
linear interpolation on each intervening interval $[1/Q_{n+1},1/Q_n]$ (with
$h_k(t)=h_k(1/Q_1)$ for $t\ge1/Q_1$). Only its values near zero matter. Write
$\mathcal H^{h_k}$ for Hausdorff measure with covering cost
$h_k(\operatorname{diam} U)$.

\begin{theorem}[Size of a critical spectral level, Layer~A in Lean]\label{thm:gauge}
On the circle $\R/\Z$, the exact level $\{x:\Lam(x)=k/4\}$ has
infinite but $\sigma$-finite $\mathcal H^{h_k}$ measure.
The strict sublevel $\{x:\Lam(x)<k/4\}$ has zero $\mathcal H^{h_k}$ measure.
In the family of gauges $g_\beta$ given by $g_\beta(1/Q_n)=\beta^{-n}$
and the same linear interpolation ($\beta>1$), the exact level has zero
$\mathcal H^{g_\beta}$ measure for $\beta>2k+1$, whereas for $1<\beta<2k+1$ its
$\mathcal H^{g_\beta}$ measure is not $\sigma$-finite.
\end{theorem}

\begin{proof}
\emph{Digits.} Use $M_j$, $u_j=Q_jx-M_j$, $A_j=4j+2$ and
$d_j=M_{j+1}-A_jM_j-M_{j-1}=A_ju_j+u_{j-1}-u_{j+1}$ as in
Section~\ref{sec:class}, with a fixed Borel rule at ties (for instance,
$M_j=\lfloor Q_jx+1/2\rfloor$). Ties occur only on the
countable set $\bigcup_{j\ge1}\{x:Q_jx\in\Z+\tfrac12\}$, which is null for every
non-atomic measure and is avoided by the Cantor construction below.
If $\Lam(x)\le k/4$, eventually $|u_j|\le(k+1)/j$, so
$d_j-4ju_j=2u_j+u_{j-1}-u_{j+1}\to0$ and
$\limsup_{j\to\infty}|d_j|=4\Lam(x)\le k$. Since $d_j\in\Z$, the integer digits
eventually satisfy $|d_j|\le k$. For the \emph{strict} sublevel $\Lam(x)<k/4$,
Corollary~\ref{thm:sub} improves this to $|d_j|\le k-1$ eventually.

\emph{Upper bounds.} Partition the non-strict sublevel $\{\Lam\le k/4\}$ on
$[0,1]$ (or $\R/\Z$) into countably many classes
$\mathcal C(N,m_{N-1},m_N)$ indexed by a common starting index $N\ge2$ and the
two starting integers $(M_{N-1},M_N)=(m_{N-1},m_N)\in\Z^2$, requiring
$|u_j|\le(k+1)/j$ and $|d_j|\le k$ for all $j\ge N$. With a fixed Borel rounding
rule these classes are Borel sets. Because $M_{j+1}=A_jM_j+M_{j-1}+d_j$ is
determined inductively from $(M_{N-1},M_N)$ by the $n-N$ digits
$d_N,\dots,d_{n-1}\in\{-k,\dots,k\}$, by level $n>N$ a class has at most
$b^{\,n-N}$ possible values of $M_n$ ($b=2k+1$). Each value of $M_n$ confines
$x$ to the interval $[M_n/Q_n-(k+1)/(nQ_n),M_n/Q_n+(k+1)/(nQ_n)]$ of diameter
$2(k+1)/(nQ_n)\le1/Q_n$ for all $n\ge2(k+1)$ (or directly inside
$[M_n/Q_n-1/(2Q_n),M_n/Q_n+1/(2Q_n)]$ of diameter $1/Q_n$ using $|u_n|\le1/2$).
Since $1/Q_n\to0$, for every $\delta>0$ and all large $n$ this gives a finite
$\delta$-cover of $\mathcal C(N,m_{N-1},m_N)$ with $\mathcal H_\delta^{h_k}$
cost at most
\[
b^{\,n-N}h_k(1/Q_n)=b^{\,n-N}b^{-n}=b^{-N}<\infty.
\]
Letting $\delta\downarrow0$ gives
$\mathcal H^{h_k}(\mathcal C(N,m_{N-1},m_N))\le b^{-N}<\infty$. For the strict
sublevel $\{\Lam<k/4\}$, the allowed tail digits satisfy
$d_j\in\{-(k-1),\dots,k-1\}$, so the count at level $n$ uses at most
$(2k-1)^{n-N}$ values of $M_n$ with $2k-1<b=2k+1$. The $\mathcal H_\delta^{h_k}$
covering cost of each class is therefore at most
\[
(2k-1)^{n-N}h_k(1/Q_n)=b^{-N}\left(\frac{2k-1}{2k+1}\right)^{n-N}\longrightarrow0
\qquad(n\to\infty),
\]
so every class is $\mathcal H^{h_k}$-null, and countable subadditivity yields
\begin{equation}\label{eq:strict-null}
 \mathcal H^{h_k}\bigl(\{x:\Lam(x)<k/4\}\bigr)=0 .
\end{equation}

\emph{Positivity.} Fix a large integer $N\ge2$ (chosen below), set
$M_{N-1}=M_N=0$, and prescribe arbitrary digits
$\omega=(d_N,d_{N+1},\dots)\in\{-k,\ldots,k\}^{\{j\ge N\}}$ via the linear
recurrence
\[
M_{j+1}=A_jM_j+M_{j-1}+d_j\qquad(j\ge N).
\]
Define the determinant sequence $W_j:=M_{j+1}Q_j-M_jQ_{j+1}$ for $j\ge N-1$, so
that $W_{N-1}=M_NQ_{N-1}-M_{N-1}Q_N=0$. Using $M_{j+1}=A_jM_j+M_{j-1}+d_j$ and
$Q_{j+1}=A_jQ_j+Q_{j-1}$ from \eqref{eq:rec}, the $A_jM_jQ_j$ terms cancel and
we obtain the first-order recurrence
\[
W_j=(A_jM_j+M_{j-1}+d_j)Q_j-M_j(A_jQ_j+Q_{j-1})
=d_jQ_j-(M_jQ_{j-1}-M_{j-1}Q_j)=d_jQ_j-W_{j-1}
\]
for all $j\ge N$. We claim by induction on $j\ge N-1$ that
\begin{equation}\label{eq:Wj-bound}
|W_j|\le 2kQ_j\qquad\text{for all }j\ge N-1.
\end{equation}
Indeed, for $j=N-1$, $|W_{N-1}|=0\le2kQ_{N-1}$. For $j\ge N$, since $Q_j\ge2Q_{j-1}$
(in fact $Q_j\ge(4j-2)Q_{j-1}\ge6Q_{j-1}$ for $j\ge2$), the inductive
hypothesis $|W_{j-1}|\le2kQ_{j-1}\le kQ_j$ and $|d_j|\le k$ yield
\[
|W_j|\le|d_j|Q_j+|W_{j-1}|\le kQ_j+2kQ_{j-1}\le kQ_j+kQ_j=2kQ_j,
\]
proving \eqref{eq:Wj-bound}. Dividing $W_j$ by $Q_jQ_{j+1}$ gives the adjacent
rational difference
\[
\frac{M_{j+1}}{Q_{j+1}}-\frac{M_j}{Q_j}=\frac{W_j}{Q_jQ_{j+1}},\qquad
\left|\frac{M_{j+1}}{Q_{j+1}}-\frac{M_j}{Q_j}\right|\le\frac{2k}{Q_{j+1}}.
\]
Because $Q_{j+2}\ge(4j+6)Q_{j+1}\ge2Q_{j+1}$ for all $j\ge1$, we have
$Q_{m+1+r}\ge2^r Q_{m+1}$ for all $m\ge N$ and $r\ge0$. Consequently, the
telescoping series converges uniformly on $\{-k,\dots,k\}^{\{j\ge N\}}$ to a
real limit
\[
x(\omega):=\lim_{j\to\infty}\frac{M_j}{Q_j}
=\frac{M_m}{Q_m}+\sum_{j=m}^\infty\frac{W_j}{Q_jQ_{j+1}}\qquad(m\ge N),
\]
with the uniform tail bound
\begin{equation}\label{eq:telescoping-tail}
\left|x(\omega)-\frac{M_m}{Q_m}\right|
\le\sum_{j=m}^\infty\frac{|W_j|}{Q_jQ_{j+1}}
\le\sum_{j=m}^\infty\frac{2k}{Q_{j+1}}
\le\frac{2k}{Q_{m+1}}\sum_{r=0}^\infty 2^{-r}
=\frac{4k}{Q_{m+1}}\qquad\text{for all }m\ge N.
\end{equation}
Multiplying \eqref{eq:telescoping-tail} by $Q_m$ gives, for every $m\ge N$,
\[
|Q_mx(\omega)-M_m|\le\frac{4kQ_m}{Q_{m+1}}\le\frac{4k}{4m+2}<\frac{k}{m}.
\]
Choose $N>8k$ (so $4k/Q_{N+1}<1/8$). Then for every $j\ge N$,
$|Q_jx(\omega)-M_j|<k/j<1/8<1/2$, which ensures that $M_j$ is the \emph{unique}
nearest integer to $Q_jx(\omega)$, that $u_j(x(\omega))=Q_jx(\omega)-M_j$ avoids
all half-integer ties, and that the true recurrence digit of $x(\omega)$ at
index $j$ is $M_{j+1}-A_jM_j-M_{j-1}=d_j$. Moreover,
$|4ju_j(x(\omega))-d_j|\le2|u_j|+|u_{j-1}|+|u_{j+1}|\to0$, so
$4\Lam(x(\omega))=\limsup_{j\to\infty}|d_j|\le k$. Hence the compact image
$K:=x(\{-k,\dots,k\}^{\{j\ge N\}})$ lies in the \emph{non-strict} sublevel
$\{x\in\R:\Lam(x)\le k/4\}$; because $M_N=0$ and $|x(\omega)|\le4k/Q_{N+1}$ by
\eqref{eq:telescoping-tail} at $m=N$, $K$ is contained in
$[-4k/Q_{N+1},4k/Q_{N+1}]$, an interval of length $8k/Q_{N+1}<1/4<1/2$ on which
the canonical projection $\R\to\R/\Z$ is an isometric embedding.

Equip $\{-k,\dots,k\}^{\{j\ge N\}}$ with the uniform Bernoulli product
probability measure assigning independent weight $b^{-1}=(2k+1)^{-1}$ to each
digit $d_j$, and let $\mu$ be its pushforward probability measure on $K$. Take
any interval $U\subset\R$ of length $h=\operatorname{diam}(U)$ with
$1/Q_{n+1}\le h<1/Q_n$ for some $n>N$. If $x,y\in K\cap U$, then for every index
$j$ with $N\le j<n$, we have $|x-y|\le h<1/Q_n$, $|Q_jx-M_j(x)|<1/8$, and
$|Q_jy-M_j(y)|<1/8$, so the triangle inequality gives
\[
|M_j(x)-M_j(y)|
\le|M_j(x)-Q_jx|+Q_j|x-y|+|Q_jy-M_j(y)|
<\frac18+\frac{Q_j}{Q_n}+\frac18=\frac14+\frac{Q_j}{Q_n}\le\frac14+\frac16<1.
\]
Since $M_j(x)$ and $M_j(y)$ are integers, $M_j(x)=M_j(y)$ for all $N\le j<n$
(and $M_{N-1}=0$ by definition). At level $n$, the same estimate gives
\[
|M_n(x)-M_n(y)|<\frac14+Q_n h<\frac14+1=\frac54<2,
\]
so for all $x\in K\cap U$ the integer $M_n(x)$ can take at most $3$ consecutive
integer values (differing by at most $1$ from $\min_{x\in K\cap U}M_n(x)$).
Because the integer sequence $M_{N-1},M_N,\dots,M_n$ uniquely determines the
first $n-N$ digits $d_N,\dots,d_{n-1}$ via $d_j=M_{j+1}-A_jM_j-M_{j-1}$, the
intersection $K\cap U$ is contained in at most $3$ cylinder sets of depth $n-N$,
each of $\mu$-measure $b^{-(n-N)}$. Since $h_k$ is nondecreasing and
$h\ge1/Q_{n+1}$, we have $h_k(h)\ge h_k(1/Q_{n+1})=b^{-(n+1)}$, and therefore
\[
\mu(U)=\mu(K\cap U)\le 3b^{-(n-N)}=3b^{N+1}b^{-(n+1)}\le 3b^{N+1}h_k(h).
\]
Summing over any countable cover $\{U_i\}$ of $K$ with
$\operatorname{diam}(U_i)<1/Q_{N+1}$ gives
$\sum_i h_k(\operatorname{diam}U_i)\ge(3b^{N+1})^{-1}\sum_i\mu(U_i)\ge(3b^{N+1})^{-1}>0$
(the mass distribution principle for $\mathcal H^{h_k}$), so
$\mathcal H^{h_k}(K)\ge(3b^{N+1})^{-1}>0$. By the isometric embedding of
$[-4k/Q_{N+1},4k/Q_{N+1}]$ into $\R/\Z$, the projection of $K$ to the circle has
the same positive $\mathcal H^{h_k}$ measure. Since $K\subseteq\{\Lam\le k/4\}$
and \eqref{eq:strict-null} shows that the strict sublevel $\{\Lam<k/4\}$ has
$\mathcal H^{h_k}$ measure zero, we conclude that
\[
 \mathcal H^{h_k}\bigl(\{x\in\R/\Z:\Lam(x)=k/4\}\bigr)
 \ \ge\ \mathcal H^{h_k}(K)-\mathcal H^{h_k}\bigl(\{\Lam<k/4\}\bigr)
 \ =\ \mathcal H^{h_k}(K)\ >\ 0 .
\]

\emph{Measurability, infinitude and $\sigma$-finiteness.} The functions
$x\mapsto j\|Q_jx\|$ are continuous on the circle $\R/\Z$. Hence $\Lam$ is Borel
measurable as their pointwise limit superior, every sublevel $\{\Lam\le c\}$ and
superlevel $\{\Lam\ge c\}=\bigcap_{m\ge1}\bigcup_{j\ge m}\{x:j\|Q_jx\|\ge c-1/m\}$
is Borel, and so is the exact level
$S_k:=\{x\in\R/\Z:\Lam(x)=k/4\}=\{\Lam\le k/4\}\cap\{\Lam\ge k/4\}$. By
Corollary~\ref{cor:infinite-level} (or Theorem~\ref{thm:finite}), $S_k$ has
Lebesgue measure zero. For any $y\in L=\Z+\Z\frac{\ee+1}{2}$, the triangle
inequality for $\|\cdot\|$ gives
$\bigl|j\|Q_j(x+y)\|-j\|Q_jx\|\bigr|\le j\|Q_jy\|\to0$ as $j\to\infty$ by
Theorem~\ref{thm:class} (the critical weight $j$ is essential here), so
$\Lam(x+y)=\Lam(x)$ and $S_k$ is invariant under translation by the subgroup
$\bar L:=L/\Z\subset\R/\Z$. Since $(\ee+1)/2$ is irrational, Kronecker's theorem
implies that $\bar L$ is dense in $\R/\Z$.

Suppose for a contradiction that $0<\mathcal H^{h_k}(S_k)<\infty$. Then the
normalized restriction
\[
\nu(B):=\frac{\mathcal H^{h_k}(S_k\cap B)}{\mathcal H^{h_k}(S_k)}
\qquad(B\subseteq\R/\Z\text{ Borel})
\]
is a Borel probability measure on $\R/\Z$ concentrated on $S_k$ ($\nu(S_k)=1$).
Because $\mathcal H^{h_k}$ is translation-invariant on $\R/\Z$ and $S_k+\ell=S_k$
for every $\ell\in\bar L$, $\nu$ is invariant under rotation by every
$\ell\in\bar L$: $\nu(B+\ell)=\nu(B)$ for all Borel $B\subseteq\R/\Z$. We show
that $\nu$ is invariant under rotation by \emph{every} $t\in\R/\Z$. Let
$f:\R/\Z\to\R$ be any continuous function. By uniform continuity on the compact
group $\R/\Z$, for any $t\in\R/\Z$ we may choose a sequence $\ell_m\in\bar L$
with $\ell_m\to t$ in $\R/\Z$, and then
$\sup_{x\in\R/\Z}|f(x-\ell_m)-f(x-t)|\to0$ as $m\to\infty$. By $\bar L$-invariance
of $\nu$ and bounded convergence,
\[
\int_{\R/\Z}f(x-t)\,d\nu(x)
=\lim_{m\to\infty}\int_{\R/\Z}f(x-\ell_m)\,d\nu(x)
=\lim_{m\to\infty}\int_{\R/\Z}f(x)\,d\nu(x)
=\int_{\R/\Z}f(x)\,d\nu(x).
\]
Thus the rotated Borel probability measure $\nu_t(B):=\nu(B+t)$ has the same
integral as $\nu$ on every continuous function $f\in C(\R/\Z)$ (or equivalently
the same Fourier coefficients $\int_{\R/\Z}\ee^{2\pi i m x}\,d\nu(x)=0$ for all
$m\in\Z\setminus\{0\}$), so $\nu_t=\nu$ for all $t\in\R/\Z$ by uniqueness of
Haar probability measure on $\R/\Z$: concretely, for every integer $m\ge1$,
partitioning $\R/\Z$ into $m$ disjoint half-open arcs of equal length $1/m$
gives mass $1/m$ to each arc, Hence every arc of rational length $p/m$ has
$\nu$-measure $p/m$, and monotonicity extends $\nu(I)=\operatorname{Leb}_{\R/\Z}(I)$
to all open arcs $I\subset\R/\Z$, which generate the Borel $\sigma$-algebra.
Thus $\nu=\operatorname{Leb}_{\R/\Z}$, contradicting
$1=\nu(S_k)=\operatorname{Leb}_{\R/\Z}(S_k)=0$. Therefore
$\mathcal H^{h_k}(S_k)=\infty$. Finally, $S_k$ has $\sigma$-finite
$\mathcal H^{h_k}$ measure because $S_k\subseteq\{\Lam\le k/4\}$, which is
covered by the countably many Borel classes $\mathcal C(N,m_{N-1},m_N)$ above,
each of finite $\mathcal H^{h_k}$ measure $\le b^{-N}<\infty$.

For the power-gauge family $g_\beta(1/Q_n)=\beta^{-n}$, on each linear
interpolation interval $[1/Q_{n+1},1/Q_n]$ both $g_\beta$ and $h_k$ are affine
with positive values at the endpoints, so the ratio $g_\beta(r)/h_k(r)$ is
monotone on $[1/Q_{n+1},1/Q_n]$ and therefore lies between the endpoint values
$(\beta/(2k+1))^{-n}$ and $(\beta/(2k+1))^{-n-1}$. Thus
$g_\beta(r)/h_k(r)\to0$ when $\beta>2k+1$ and $h_k(r)/g_\beta(r)\to0$ when
$1<\beta<2k+1$ as $r\downarrow0$, which are special cases of
Corollary~\ref{cor:exact-gauge}(ii)--(iii).

For comparison, Cabrelli, Mendivil, Molter and Shonkwiler
\cite[\S5, Theorem~3]{CMMS} construct, for a Cantor set whose gap
sequence is nonincreasing, a gauge $h$ for which the set is an $h$-set
($1/32\le\mathcal H^h\le1$, a finite positive measure), whereas Moreira, Erazo
and Angelini \cite{MEA} show that the classical Lagrange level at value $3$ is
never $\sigma$-finite with positive measure for any gauge. Here the gauge comes
from the recurrence's own scale, and the entire exact level has \emph{infinite}
but $\sigma$-finite $\mathcal H^{h_k}$ measure.
\end{proof}

\paragraph{Lean status.}
The null part \eqref{eq:strict-null} is kernel-checked at Layer~A in
\texttt{EpiFormal/\allowbreak GaugeNull.lean} as
\texttt{GaugeNull.\allowbreak mkMetric\_\allowbreak strictLevel\_\allowbreak null}
for any nondecreasing function $m:\R_{\ge0\infty}\to\R_{\ge0\infty}$ satisfying
$m(1/Q_n)\le C(2k+1)^{-n}$ ($n\ge1$), without requiring continuity of $m$; the
same module kernel-checks Corollary~\ref{cor:exact-gauge}(ii) on $\R$ as
\texttt{GaugeNull.\allowbreak mkMetric\_\allowbreak levelSet\_\allowbreak null\_\allowbreak of\_\allowbreak subgauge}.
In \texttt{EpiFormal/\allowbreak GaugeExact.lean}, $L$-translation invariance
(\texttt{exactLevel\_\allowbreak add\_\allowbreak latticeL}), the countable
finite-measure covering
(\texttt{exactLevel\_\allowbreak countable\_\allowbreak cover\_\allowbreak finite\_\allowbreak mkMetric}),
and---for arbitrary monotone continuous $m$ satisfying matching two-sided
bounds---the Cantor-tree lower bound
(\texttt{mkMetric\_\allowbreak cantorTree\_\allowbreak pos}), exact-level
positivity (\texttt{mkMetric\_\allowbreak exactLevel\_\allowbreak pos}),
exact-level infinitude on $\R$
(\texttt{mkMetric\_\allowbreak exactLevel\_\allowbreak eq\_\allowbreak top}),
and the non-$\sigma$-finiteness criterion
(\texttt{not\_\allowbreak countable\_\allowbreak union\_\allowbreak finite\_\allowbreak of\_\allowbreak ratio\_\allowbreak tendsto\_\allowbreak zero})
are likewise kernel-checked.

The external continuity hypothesis, the explicit piecewise-linear gauge
construction, and the native quotient-metric circle statements on
\texttt{UnitAddCircle} ($\R/\Z$) are completely discharged in
\texttt{EpiFormal/\allowbreak GaugeCircle.lean}:
the gauges $h_k$ (\texttt{critGauge k}), $g_2$ (\texttt{binaryGauge}), and
$g_\beta$ (\texttt{interpGauge \ensuremath{\beta}}) are constructed explicitly
with kernel-checked monotonicity (\texttt{critGauge\_\allowbreak mono}),
continuity on $\R_{\ge0\infty}$ including $0^+$ and $\infty$
(\texttt{critGauge\_\allowbreak continuous}), and exact endpoint and segment
values (\texttt{critGauge\_\allowbreak at\_\allowbreak inv\_\allowbreak Qsel},
\texttt{critGauge\_\allowbreak eq\_\allowbreak segVal}).
This yields the assumption-free instantiations of Theorem~\ref{thm:gauge},
Corollary~\ref{cor:exact-gauge}(i)--(iii), and Remark~\ref{rem:gauge-scale} both
on $\R$
(\texttt{mkMetric\_\allowbreak critGauge\_\allowbreak strictLevel\_\allowbreak null},
\texttt{mkMetric\_\allowbreak critGauge\_\allowbreak exactLevel\_\allowbreak pos},
\texttt{mkMetric\_\allowbreak critGauge\_\allowbreak exactLevel\_\allowbreak eq\_\allowbreak top},
\texttt{critGauge\_\allowbreak exactLevel\_\allowbreak sigma\_\allowbreak finite\_\allowbreak cover},
\texttt{mkMetric\_\allowbreak critGauge\_\allowbreak exactLevel\_\allowbreak inter\_\allowbreak Ioo\_\allowbreak eq\_\allowbreak top},
\texttt{mkMetric\_\allowbreak powerGauge\_\allowbreak exactLevel\_\allowbreak null\_\allowbreak of\_\allowbreak gt},
\texttt{mkMetric\_\allowbreak powerGauge\_\allowbreak exactLevel\_\allowbreak not\_\allowbreak sigma\_\allowbreak finite\_\allowbreak of\_\allowbreak lt},
and \texttt{critGauge\_\allowbreak adjacent\_\allowbreak level\_\allowbreak separation})
and natively on \texttt{UnitAddCircle} ($\R/\Z$)
(\texttt{mkMetric\_\allowbreak critGauge\_\allowbreak circleStrictLevel\_\allowbreak null},
\texttt{mkMetric\_\allowbreak circleLevelSet\_\allowbreak null\_\allowbreak of\_\allowbreak subgauge},
\texttt{mkMetric\_\allowbreak critGauge\_\allowbreak circleExactLevel\_\allowbreak inter\_\allowbreak open\_\allowbreak eq\_\allowbreak top},
\texttt{mkMetric\_\allowbreak critGauge\_\allowbreak circleExactLevel\_\allowbreak eq\_\allowbreak top},
\texttt{critGauge\_\allowbreak circleExactLevel\_\allowbreak sigma\_\allowbreak finite\_\allowbreak cover},
\texttt{mkMetric\_\allowbreak critGauge\_\allowbreak circleExactLevel\_\allowbreak not\_\allowbreak sigma\_\allowbreak finite\_\allowbreak of\_\allowbreak ratio},
\texttt{mkMetric\_\allowbreak powerGauge\_\allowbreak circleExactLevel\_\allowbreak null\_\allowbreak of\_\allowbreak gt},
and
\texttt{mkMetric\_\allowbreak powerGauge\_\allowbreak circleExactLevel\_\allowbreak not\_\allowbreak sigma\_\allowbreak finite\_\allowbreak of\_\allowbreak lt}),
proved constructively by splitting the Cantor tree into depth-$M$ cylinders,
packing $\Omega\bigl((b-a)Q_{J+M+1}\bigr)$ disjoint $L$-translates using
$\delta_j/2\in L$, and transferring across the local isometry
$\R\to\R/\Z$ on intervals of length $\le1/2$. For the binary exact-limit
regime of Theorem~\ref{thm:trichotomy}(ii), the binary Cantor-tree lower bound,
global and open-set $\mathcal H^{g_2}$-infinitude, $\sigma$-finite covering,
$\mathcal H^{h_k}$-nullity, and the $g_\beta$ threshold dichotomy ($\beta>2$ vs.\
$1<\beta<2$) are likewise kernel-checked on both $\R$ and \texttt{UnitAddCircle}
(\texttt{mkMetric\_\allowbreak binaryGauge\_\allowbreak unsigned\allowbreak ExactLimit\_\allowbreak pos},
\texttt{mkMetric\_\allowbreak binaryGauge\_\allowbreak unsigned\allowbreak ExactLimit\_\allowbreak inter\_\allowbreak Ioo\_\allowbreak eq\_\allowbreak top},
\texttt{mkMetric\_\allowbreak binaryGauge\_\allowbreak circle\allowbreak Unsigned\allowbreak ExactLimit\_\allowbreak inter\_\allowbreak open\_\allowbreak eq\_\allowbreak top},
\texttt{mkMetric\_\allowbreak binaryGauge\_\allowbreak circle\allowbreak Unsigned\allowbreak ExactLimit\_\allowbreak eq\_\allowbreak top},
\texttt{binaryGauge\_\allowbreak circle\allowbreak Unsigned\allowbreak ExactLimit\_\allowbreak sigma\_\allowbreak finite\_\allowbreak cover},
\texttt{critGauge\_\allowbreak circle\allowbreak Unsigned\allowbreak ExactLimit\_\allowbreak null},
\texttt{mkMetric\_\allowbreak powerGauge\_\allowbreak circle\allowbreak Unsigned\allowbreak ExactLimit\_\allowbreak null\_\allowbreak of\_\allowbreak gt},
and
\texttt{mkMetric\_\allowbreak powerGauge\_\allowbreak circle\allowbreak Unsigned\allowbreak ExactLimit\_\allowbreak not\_\allowbreak sigma\_\allowbreak finite\_\allowbreak of\_\allowbreak lt}).
All 53 audited declarations in \texttt{GaugeCircle.lean} report only
\texttt{propext}, \texttt{Classical.choice}, and \texttt{Quot.sound}, with zero
\texttt{sorryAx}.

The gauge $h_k$ is exact for the level $k/4$ in the following sense.
Here a gauge is a nondecreasing function $g:[0,\infty)\to[0,\infty)$
with $g(0)=0$, and $\mathcal H^g$ is the Hausdorff measure on the circle
with covering cost $g(\operatorname{diam}U)$.

\begin{corollary}[Exact critical gauge]\label{cor:exact-gauge}
Let $k\ge1$ and $S_k=\{x\in\R/\Z:\Lam(x)=k/4\}$.
\begin{enumerate}[label=(\roman*)]
\item For every nonempty open arc $I\subseteq\R/\Z$,
$\mathcal H^{h_k}(S_k\cap I)=\infty$.
\item If $g$ is a gauge with $g(1/Q_n)/h_k(1/Q_n)\to0$ as $n\to\infty$, then
$\mathcal H^{g}(\{\Lam\le k/4\})=0$; in particular $\mathcal H^g(S_k)=0$.
(Lean on $\R$ and $\R/\Z$: \texttt{GaugeNull.\allowbreak mkMetric\_\allowbreak levelSet\_\allowbreak null\_\allowbreak of\_\allowbreak subgauge} and \texttt{GaugeCircle.\allowbreak mkMetric\_\allowbreak circleLevelSet\_\allowbreak null\_\allowbreak of\_\allowbreak subgauge}.)
\item If $g$ is a gauge with $h_k(t)/g(t)\to0$ as $t\to0^+$, then $S_k$ is
not $\sigma$-finite for $\mathcal H^g$; in particular
$\mathcal H^g(S_k)=\infty$.
\end{enumerate}
\end{corollary}

\begin{proof}
(i) Hausdorff measures on the circle are invariant under rotations, and
$S_k$ is invariant under rotation by the elements of $L$ (proof of
Theorem~\ref{thm:gauge}). Suppose $\mathcal H^{h_k}(S_k\cap I)<\infty$ for
some nonempty open arc $I$. Since $L$ is dense modulo one, the arcs $I+\ell$
($\ell\in L$) cover the circle, and by compactness finitely many of them,
say $I+\ell_1,\ldots,I+\ell_N$, already do. Then
$\mathcal H^{h_k}(S_k)\le\sum_{i=1}^N\mathcal H^{h_k}(S_k\cap(I+\ell_i))
=N\,\mathcal H^{h_k}(S_k\cap I)<\infty$, contradicting
Theorem~\ref{thm:gauge}.

(ii) In the proof of Theorem~\ref{thm:gauge}, the non-strict sublevel
$\{\Lam\le k/4\}$ is a countable union of classes. For a class with
starting index $N$ and every large $n$, at most $b^{\,n-N}$ intervals of
diameter at most $1/Q_n$ cover it, where $b=2k+1$. Since $g$ is
nondecreasing, the $g$-cost of this cover is at most
\[
 b^{\,n-N}g(1/Q_n)=b^{-N}\,\frac{g(1/Q_n)}{h_k(1/Q_n)}\longrightarrow0 .
\]
So every class is $\mathcal H^g$-null, and so is their countable union.

(iii) Fix $\varepsilon>0$ and choose $\delta>0$ with
$h_k(t)\le\varepsilon g(t)$ for $0\le t<\delta$. Every $\delta$-cover then
satisfies $\sum h_k(\operatorname{diam}U)\le
\varepsilon\sum g(\operatorname{diam}U)$, so
$\mathcal H^{h_k}(E)\le\varepsilon\,\mathcal H^g(E)$ for every set $E$.
Hence $\mathcal H^g(E)<\infty$ implies $\mathcal H^{h_k}(E)=0$. If $S_k$
were a countable union of sets $E_i$ with $\mathcal H^g(E_i)<\infty$, we
would get $\mathcal H^{h_k}(S_k)=0$, contradicting Theorem~\ref{thm:gauge}.
\end{proof}

\begin{corollary}[Arbitrary gauges at the exact and unsigned levels]\label{cor:liminf-gauges}
Let $g$ be a Hausdorff gauge, $r_n=1/Q_n$, and
\[
\rho_b(g)=\liminf_{n\to\infty} b^n g(r_n).
\]
For $S_k=\{\Lam=k/4\}$ take $b=2k+1$, and for
$T_k=\{x:j\|Q_jx\|\to k/4\}$ take $b=2$, where $k\ge1$ and both sets
are taken on the circle. In each case the named set is
$\mathcal H^g$-null if $\rho_b(g)=0$. If $0<\rho_b(g)<\infty$, it is
$\sigma$-finite and its intersection with every nonempty open set has
infinite $\mathcal H^g$ measure. If $\rho_b(g)=\infty$, its intersection
with every nonempty open set is non-$\sigma$-finite. For $b=2k+1$ the
zero assertion applies to the full sublevel $\{\Lam\le k/4\}$, and
$\rho_b(g)<\infty$ implies that $\{\Lam<k/4\}$ is $\mathcal H^g$-null.
\end{corollary}

\begin{proof}
The fixed initial-pair classes of Theorems~\ref{thm:gauge}
and~\ref{thm:trichotomy}(ii) have covers by at most $b^{n-N}$ intervals
of diameter $r_n$ for a fixed starting level $N$ and all large $n$.
Their gauge cost is bounded by $b^{-N}b^n g(r_n)$. Taking the lower
limit along permissible scales gives zero or finite measure for each
class, according as $\rho_b(g)$ is zero or finite, and the classes are
countable. For the strict sublevel the same argument uses $2k-1$
instead of $2k+1$. Bounded cost along a subsequence at base $2k+1$
then tends to zero along that subsequence at base $2k-1$.

Write $h_b(r_n)=b^{-n}$ with linear interpolation. If
$r_{n+1}\le r\le r_n$, then
\[
 h_b(r)\le b^{-n},\qquad g(r)\ge g(r_{n+1}).
\]
For any $0<A<\rho_b(g)$, eventually $b^{n+1}g(r_{n+1})\ge A$,
so $h_b(r)\le(b/A)g(r)$ for all sufficiently small $r$. The critical
reference measure is infinite on the intersection with every nonempty
open set, by Corollary~\ref{cor:exact-gauge}(i) and
Theorem~\ref{thm:trichotomy}(ii), respectively. The comparison gives the
same local infinitude for $g$. When $\rho_b(g)=\infty$, $A$ can be
arbitrarily large, so any set of finite $\mathcal H^g$ measure is
$\mathcal H^{h_b}$-null. A countable union of such sets cannot cover
a nonempty open intersection of the named set.
\end{proof}

\begin{remark}[Oscillating normalized costs]
Let $n_j=2^j$ and prescribe $g(r_n)=b^{-n_j}$ for
$n_j\le n<n_{j+1}$, interpolating linearly in $r$. The lower and upper
limits of $b^ng(r_n)$ are $1$ and $\infty$. Prescribing
$g(r_n)=b^{-(n_{j+1}-1)}$ instead gives lower and upper limits $0$ and $1$.
Both prescriptions are nondecreasing as functions of $r$ and tend to zero
at the origin. They can be made strictly increasing near zero by adding
$h_b^2$, which adds $b^{-n}$ to the normalized costs at the nodes.
\end{remark}

\paragraph{Lean verification.}
The accompanying \texttt{GaugeLiminfComplete.lean} verifies
Corollary~\ref{cor:liminf-gauges} on \texttt{UnitAddCircle}, including
the strict-sublevel assertion and the non-$\sigma$-finite conclusion on
each nonempty open set. Its three principal declarations are
\texttt{circleExactLevel\_\allowbreak liminf\_\allowbreak trichotomy},
\texttt{circleUnsigned\allowbreak ExactLimit\_\allowbreak liminf\_\allowbreak trichotomy}, and
\texttt{mkMetric\_\allowbreak circleStrictLevel\_\allowbreak null\_\allowbreak of\_\allowbreak liminf\_\allowbreak lt\_\allowbreak top}.
All eighteen audited declarations use only the three base axioms.

\begin{remark}[Level separation and interpolation-free logarithmic scale]\label{rem:gauge-scale}
The gauges $h_k$ strictly separate the critical levels. Indeed, for $k\ge2$,
$h_k(1/Q_n)/h_{k-1}(1/Q_n)=\bigl((2k-1)/(2k+1)\bigr)^n\to0$, and on each
interval $[1/Q_{n+1},1/Q_n]$ both gauges are linear, so their ratio lies
between its values at the endpoints. Hence $h_k(t)/h_{k-1}(t)\to0$ as
$t\to0^+$. By Corollary~\ref{cor:exact-gauge}(iii), the level $S_k$ is not
$\sigma$-finite for $\mathcal H^{h_{k-1}}$, whereas $S_{k-1}$ has infinite
but $\sigma$-finite $\mathcal H^{h_{k-1}}$ measure. By
Corollary~\ref{cor:exact-gauge}(ii), the level $S_{k-1}$ is
$\mathcal H^{h_k}$-null, consistent with \eqref{eq:strict-null}.

Moreover, the scale of the critical gauge $h_k$ can be read without reference
to the linear interpolation. For $1/Q_{n+1}\le r\le1/Q_n$ the quantity
$-\log h_k(r)$ lies between $n\log(2k+1)$ and $(n+1)\log(2k+1)$.
By \eqref{eq:rec} both endpoints for $\log(1/r)$ are asymptotic to
$n\log n$, giving
\[
 -\log h_k(r)\sim \log(2k+1)\,
             \frac{\log(1/r)}{\log\log(1/r)}\qquad(r\downarrow0).
\]
In particular, any gauge $g$ whose corresponding logarithmic rate exists and is
strictly greater than $\log(2k+1)$ satisfies $g(1/Q_n)/h_k(1/Q_n)\to0$ and
therefore gives zero measure to $S_k$, while any gauge with a strictly smaller
logarithmic rate satisfies $h_k(t)/g(t)\to0$ and gives a non-$\sigma$-finite
measure on $S_k$.
\end{remark}

\subsection{The three-way splitting at each positive spectral level}

While every subcritical weight $0\le\alpha<1$ assigns the same Hausdorff
dimension $1-\alpha$ to the sublevel set $\{\Lam_\alpha\le c\}$ and to the
exact limit set $\{j^\alpha\|Q_jx\|\to c\}$ (Theorem~\ref{conj:weight}), the
critical weight $\alpha=1$ exhibits a dramatic three-way splitting at each
positive quantized level $c=k/4$ ($k\in\Z_{\ge1}$), governed by the allowed
tail alphabet of the integer recurrence digits $d_j(x)$.

\begin{theorem}[Trichotomy at positive spectral levels, Layer~A in Lean]\label{thm:trichotomy}
Fix an integer $k\ge1$ and let $u_j(x)=Q_jx-\mathrm{round}(Q_jx)$. At the
spectral level $c=k/4$, the sublevel set, the unsigned exact limit set, and
the signed exact limit sets split into three distinct regimes:
\begin{enumerate}[label=(\roman*),leftmargin=2em]
\item \textbf{Limsup level (alphabet of size $2k+1$, gauge base $2k+1$):}
$x\in E_{k/4}=\{\Lam\le k/4\}$ if and only if $u_j(x)\to0$ and eventually
$d_j(x)\in\{-k,-k+1,\dots,k\}$. On $\R/\Z$, both $E_{k/4}$ and the exact
limsup level $\{\Lam=k/4\}$ have infinite, $\sigma$-finite $\mathcal H^{h_k}$
measure for the exact critical gauge $h_k(1/Q_n)=(2k+1)^{-n}$
(Theorem~\ref{thm:gauge} and Corollary~\ref{cor:exact-gauge}).
\item \textbf{Unsigned exact limit set (universal binary alphabet of size $2$, gauge base $2$):}
Define $E_{\mathrm{lim}}(k/4):=\{x\in\R:\lim_{j\to\infty}j\|Q_jx\|=k/4\}$.
Then $x\in E_{\mathrm{lim}}(k/4)$ if and only if $u_j(x)\to0$ and eventually
$|d_j(x)|=k$, i.e.\ $d_j(x)\in\{-k,+k\}$ (exactly two branches per step,
\emph{independently} of $k\ge1$). Consequently, in the interpolated gauge
family $g_\beta(1/Q_n)=\beta^{-n}$, the exact critical gauge for
$E_{\mathrm{lim}}(k/4)$ is $h_{\mathrm{lim}}:=g_2$ ($h_{\mathrm{lim}}(1/Q_n)=2^{-n}$)
for \emph{every} $k\ge1$: on $\R/\Z$, $E_{\mathrm{lim}}(k/4)$ has infinite,
$\sigma$-finite $\mathcal H^{g_2}$ measure (and infinite $\mathcal H^{g_2}$
measure in every nonempty open arc, whereas
$\mathcal H^{h_k}(E_{\mathrm{lim}}(k/4))=0$ since $2<2k+1$), and
$E_{\mathrm{lim}}(c)=\emptyset$ for all $c\notin\frac14\Z_{\ge0}$.
\item \textbf{Signed exact limit sets (single tail branch, single countable $L$-coset):}
For each sign $\pm$, define
$E_{\mathrm{lim}}^\pm(k/4):=\{x\in\R:\lim_{j\to\infty}j\,u_j(x)=\pm k/4\}$.
Then $x\in E_{\mathrm{lim}}^\pm(k/4)$ if and only if $u_j(x)\to0$ and
eventually $d_j(x)=\pm k$. For any $x,y\in E_{\mathrm{lim}}^\pm(k/4)$, their
difference satisfies $x-y\in L=\Z+\Z\frac{\ee+1}{2}$. Hence each signed limit
set is a single countable coset of $L$:
\[
E_{\mathrm{lim}}^+(k/4)=x_k^+ + L,\qquad
E_{\mathrm{lim}}^-(k/4)=x_k^- + L\quad(x_k^-=-x_k^+),
\]
and therefore is countable and has zero measure for every non-atomic Hausdorff
gauge.
\end{enumerate}
(Lean declarations in \texttt{Full\allowbreak Spectrum\allowbreak Discrete.lean},
\texttt{Gauge\allowbreak Circle.lean}, and
\texttt{Trichotomy\allowbreak And\allowbreak Poly\allowbreak Extension.lean}:
the binary digit characterization \texttt{unsigned\_\allowbreak exact\_\allowbreak limit\_\allowbreak iff\_\allowbreak defect\_\allowbreak binary},
the constant digit characterization \texttt{signed\_\allowbreak exact\_\allowbreak limit\_\allowbreak iff\_\allowbreak defect\_\allowbreak const},
the $L$-coset difference theorem \texttt{sub\_\allowbreak mem\_\allowbreak latticeL\_\allowbreak of\_\allowbreak same\_\allowbreak signed\_\allowbreak limit},
countability \texttt{signed\_\allowbreak exact\_\allowbreak limit\_\allowbreak set\_\allowbreak countable},
zero dimension \texttt{dimH\_\allowbreak signed\_\allowbreak exact\_\allowbreak limit\_\allowbreak set\_\allowbreak zero},
the native $\R$ and \texttt{UnitAddCircle} binary-gauge theorems for~(ii) in \texttt{GaugeCircle.lean},
and the full signed-coset collapse for~(iii) via
\texttt{xPlus\_\allowbreak mem\_\allowbreak signed\allowbreak Exact\allowbreak LimitSet},
\texttt{signed\allowbreak Exact\allowbreak LimitSet\_\allowbreak pos\_\allowbreak eq\_\allowbreak coset},
\texttt{signed\allowbreak Exact\allowbreak LimitSet\_\allowbreak neg\_\allowbreak eq\_\allowbreak coset}, and
\texttt{trichotomy\_\allowbreak part\_\allowbreak iii\_\allowbreak complete}.)
\end{theorem}

\begin{proof}
Statement~(i) is immediate from Theorem~\ref{thm:full-spectrum},
Corollary~\ref{thm:sub}, Theorem~\ref{thm:gauge}, and
Corollary~\ref{cor:exact-gauge}.

For~(ii), if $j\|Q_jx\|\to k/4$, then $u_j(x)=O(1/j)\to0$, and
\eqref{eq:digit-asymp} gives $|d_j(x)|\to k$. Since $|d_j(x)|\in\Z_{\ge0}$, we
must have $|d_j(x)|=k$, i.e.\ $d_j(x)\in\{-k,+k\}$, for all sufficiently large
$j$. Conversely, if $u_j(x)\to0$ and eventually $d_j(x)\in\{-k,+k\}$, then
$|d_j(x)|=k$ eventually, so \eqref{eq:digit-asymp} implies $4j\|Q_jx\|\to k$,
that is, $x\in E_{\mathrm{lim}}(k/4)$ (and if $c\notin\frac14\Z_{\ge0}$, no
integer sequence $|d_j|$ can converge to $4c$, so $E_{\mathrm{lim}}(c)=\emptyset$).
We now verify each metric assertion for $E_{\mathrm{lim}}(k/4)$ on $\R/\Z$ under
the binary gauge $g_2(1/Q_n)=2^{-n}$:
\begin{itemize}[leftmargin=1.8em]
\item \emph{$\sigma$-finiteness under $\mathcal H^{g_2}$ and nullity under $\mathcal H^{h_k}$:}
Partition $E_{\mathrm{lim}}(k/4)$ into countably many Borel classes
$\mathcal C_{\mathrm{bin}}(N,m_{N-1},m_N)$ indexed by a starting index $N\ge2$
and starting pair $(M_{N-1},M_N)=(m_{N-1},m_N)\in\Z^2$, requiring
$|u_j|\le(k+1)/j$ and $d_j\in\{-k,+k\}$ for all $j\ge N$. Since at each step
$j\ge N$ there are only $2$ choices for $d_j\in\{-k,+k\}$, at level $n>N$ a
class has at most $2^{n-N}$ possible values of $M_n$, each confining $x$ to an
interval of diameter $2(k+1)/(nQ_n)\le1/Q_n$ for all $n\ge2(k+1)$. Thus the
$\mathcal H_\delta^{g_2}$ covering cost of $\mathcal C_{\mathrm{bin}}(N,m_{N-1},m_N)$
is at most $2^{n-N}g_2(1/Q_n)=2^{-N}<\infty$, proving that
$E_{\mathrm{lim}}(k/4)$ is $\sigma$-finite for $\mathcal H^{g_2}$. For the
larger-base gauge $h_k=g_{2k+1}$, since $2<2k+1$ for every $k\ge1$, the
$\mathcal H_\delta^{h_k}$ covering cost of each class is at most
\[
2^{n-N}h_k(1/Q_n)=(2k+1)^{-N}\left(\frac{2}{2k+1}\right)^{n-N}\longrightarrow0
\qquad(n\to\infty),
\]
so every class is $\mathcal H^{h_k}$-null and
$\mathcal H^{h_k}(E_{\mathrm{lim}}(k/4))=0$. More generally, any gauge $g$ with
$g(1/Q_n)/g_2(1/Q_n)=2^ng(1/Q_n)\to0$ gives
$\mathcal H^g(E_{\mathrm{lim}}(k/4))=0$.
\item \emph{Positivity under $\mathcal H^{g_2}$:}
Fix $N>8k$, set $M_{N-1}=M_N=0$, and restrict the digit sequences in the
\emph{Positivity} step of Theorem~\ref{thm:gauge} to the binary space
$\{-k,+k\}^{\{j\ge N\}}$. By \eqref{eq:telescoping-tail}, every resulting point
$x(\omega)=\lim_{j\to\infty}M_j/Q_j$ satisfies $|Q_jx(\omega)-M_j|<k/j<1/8$ and
has true recurrence digits $d_j(x(\omega))=d_j\in\{-k,+k\}$ for all $j\ge N$, so
$|d_j(x(\omega))|=k$ and $j\|Q_jx(\omega)\|\to k/4$. Hence the compact binary
Cantor image $K_{\mathrm{bin}}:=x(\{-k,+k\}^{\{j\ge N\}})$ lies entirely inside
$E_{\mathrm{lim}}(k/4)\cap[-4k/Q_{N+1},4k/Q_{N+1}]$ (no strict sublevel needs to
be subtracted). Equip $\{-k,+k\}^{\{j\ge N\}}$ with the uniform Bernoulli
product measure assigning weight $1/2$ to $-k$ and $1/2$ to $+k$, and let
$\mu_{\mathrm{bin}}$ be its pushforward to $K_{\mathrm{bin}}$. As in the proof
of Theorem~\ref{thm:gauge}, any interval $U$ of length
$1/Q_{n+1}\le h<1/Q_n$ ($n>N$) determines a unique $M_j$ for each $N\le j<n$ and
at most $3$ values of $M_n$, so $K_{\mathrm{bin}}\cap U$ meets at most $3$ binary
cylinders of depth $n-N$, each of $\mu_{\mathrm{bin}}$-mass $2^{-(n-N)}$. Thus
\[
\mu_{\mathrm{bin}}(U)\le 3\cdot 2^{-(n-N)}=3\cdot 2^{N+1}\cdot 2^{-(n+1)}
\le 3\cdot 2^{N+1}g_2(h),
\]
and the mass distribution principle gives
$\mathcal H^{g_2}(E_{\mathrm{lim}}(k/4))\ge\mathcal H^{g_2}(K_{\mathrm{bin}})\ge(3\cdot 2^{N+1})^{-1}>0$.
\item \emph{Infinitude on every nonempty open arc and non-$\sigma$-finiteness for larger gauges:}
For any $y\in L$, $j\|Q_jy\|\to0$ by Theorem~\ref{thm:class}, so
$\bigl|j\|Q_j(x+y)\|-j\|Q_jx\|\bigr|\le j\|Q_jy\|\to0$; hence
$E_{\mathrm{lim}}(k/4)$ is Borel, Lebesgue-null (since
$E_{\mathrm{lim}}(k/4)\subseteq E_{k/4}$), and invariant under translation by
the dense subgroup $\bar L=L/\Z$ on $\R/\Z$. The Haar-uniqueness argument of
Theorem~\ref{thm:gauge} therefore forces
$\mathcal H^{g_2}(E_{\mathrm{lim}}(k/4))=\infty$ on $\R/\Z$, and the compactness
covering argument of Corollary~\ref{cor:exact-gauge}(i),(iii) with $g_2$ in
place of $h_k$ shows that $\mathcal H^{g_2}(E_{\mathrm{lim}}(k/4)\cap I)=\infty$
for every nonempty open arc $I\subseteq\R/\Z$ and that $E_{\mathrm{lim}}(k/4)$
is not $\sigma$-finite for $\mathcal H^g$ whenever $g_2(t)/g(t)\to0$ as
$t\to0^+$ (in particular for $g_\beta$ with $1<\beta<2$).
\end{itemize}

For~(iii), the signed asymptotic $d_j(x)-4ju_j(x)\to0$ in
\eqref{eq:digit-asymp} shows that $ju_j(x)\to\pm k/4$ if and only if
$u_j(x)\to0$ and the integer $d_j(x)$ is eventually constant and equal to
$\pm k$. The existence of a witness $x_k^+\in E_{\mathrm{lim}}^+(k/4)$ (and
$x_k^-=-x_k^+\in E_{\mathrm{lim}}^-(k/4)$) is given by the constant-digit
construction $d_j\equiv k$ in Theorem~\ref{thm:gauge}. If
$x,y\in E_{\mathrm{lim}}^\pm(k/4)$, then $j(u_j(x)-u_j(y))\to0$; since
$Q_j(x-y)=\bigl(M_j(x)-M_j(y)\bigr)+\bigl(u_j(x)-u_j(y)\bigr)$ with
$M_j(x)-M_j(y)\in\Z$, we have $\|Q_j(x-y)\|\le|u_j(x)-u_j(y)|$, hence
$j\|Q_j(x-y)\|\to0$. By Theorem~\ref{thm:class}, $x-y\in L$. Conversely, for
any $z\in L$, $j\|Q_jz\|\to0$ gives $u_j(z)=o(1/j)\to0$; for
$x\in E_{\mathrm{lim}}^\pm(k/4)$, since $u_j(x)=O(1/j)\to0$, we have
$|u_j(x)+u_j(z)|<1/2$ for all large $j$, so $Q_j(x+z)$ rounds to
$M_j(x)+M_j(z)$ and $u_j(x+z)=u_j(x)+u_j(z)$, whence
$j\,u_j(x+z)=j\,u_j(x)+j\,u_j(z)\to\pm k/4$. Thus $x+z\in E_{\mathrm{lim}}^\pm(k/4)$
and $E_{\mathrm{lim}}^\pm(k/4)=x_k^\pm+L$.
\end{proof}

\paragraph{Lean status for Theorem~\ref{thm:trichotomy}.}
All three parts~(i), (ii), and~(iii) of Theorem~\ref{thm:trichotomy} are now
kernel-checked at Layer~A (using only the three Lean base axioms
\texttt{propext}, \texttt{Classical.choice}, and \texttt{Quot.sound}, with zero
\texttt{sorryAx}). Specifically, the digit characterizations, signed-limit
$L$-difference, countability, and dimension-zero lemmas are proved in
\texttt{EpiFormal/\allowbreak Full\allowbreak Spectrum\allowbreak Discrete.lean};
all metric gauge statements for the limsup regime~(i) and the binary exact-limit
regime $E_{\mathrm{lim}}(k/4)$ in~(ii)---binary Cantor-tree positivity, global
and nonempty-open-subset $\mathcal H^{g_2}$-infinitude, $\sigma$-finite
covering, $\mathcal H^{h_k}$-nullity, and the $g_\beta$ threshold dichotomy
($\beta>2$ vs.\ $1<\beta<2$)---are proved natively on both $\R$ and
\texttt{UnitAddCircle} ($\R/\Z$) in
\texttt{EpiFormal/\allowbreak GaugeCircle.lean} (building on
\texttt{GaugeNull.lean} and \texttt{GaugeExact.lean}); and the full
signed-coset collapse $E_{\mathrm{lim}}^\pm(k/4)=x_k^\pm+L$ in part~(iii) is
proved in
\texttt{EpiFormal/\allowbreak TrichotomyAndPolyExtension.lean} across eight
audited declarations (\texttt{xPlus\_\allowbreak mem\_\allowbreak signedExactLimitSet},
\texttt{neg\_\allowbreak xPlus\_\allowbreak mem\_\allowbreak signedExactLimitSet},
\texttt{add\_\allowbreak latticeL\_\allowbreak mem\_\allowbreak signedExactLimitSet},
\texttt{signedExactLimitSet\_\allowbreak pos\_\allowbreak eq\_\allowbreak coset},
\texttt{signedExactLimitSet\_\allowbreak neg\_\allowbreak eq\_\allowbreak coset},
\texttt{signedExactLimitSet\_\allowbreak zero\_\allowbreak eq\_\allowbreak latticeL},
\texttt{signedExactLimitSet\_\allowbreak empty\_\allowbreak of\_\allowbreak not\_\allowbreak quarter\_\allowbreak int},
and \texttt{trichotomy\_\allowbreak part\_\allowbreak iii\_\allowbreak complete}).

\subsection{Finite-state constraints on defect tails}\label{sec:finite-state-tail}

Let $G$ be a finite directed graph with edge labels in
$\{-k,\ldots,k\}$, with distinct labels on the outgoing edges at each
vertex. Mark the edges labelled $k$ or $-k$. Let $S_k(G)\subset S_k$ be
the points whose eventual defect sequence labels an infinite path in $G$.
A strongly connected component is called marked if it contains an internal
marked edge. Write $\lambda_C$ for the Perron root of the adjacency matrix
of a marked component $C$, counting edge multiplicities.

\begin{theorem}[Finite-state defect tails]\label{thm:finite-state-tail}
If $G$ has no marked component, then $S_k(G)=\varnothing$.
Otherwise put $\lambda_*=\max_C\lambda_C$, where $C$ ranges over
the marked components. If $\lambda_*=1$, then $S_k(G)$ is countable.
If $\lambda_*>1$, it satisfies the three measure conclusions of
Corollary~\ref{cor:liminf-gauges}, with $S_k(G)$ in place of the named
set and with $b=\lambda_*$. In the zero case the null conclusion refers
to $S_k(G)$.
\end{theorem}

\begin{proof}
First consider a strongly connected marked graph with Perron root
$\lambda>1$. Take a positive right Perron vector $w$. The sum of
terminal weights over paths of length $m$ from vertex $v$ equals
$\lambda^m w_v$. Since the coordinates of $w$ are bounded above and
away from zero, the path count is comparable to $\lambda^m$, uniformly
over the initial vertex. The fixed initial-pair covering argument
therefore gives at most $C\lambda^{n-N}$ intervals of diameter $r_n$
for each fixed starting position, pair and vertex. This gives the
zero and finite-cover conclusions when the normalized lower limit is
zero or finite, respectively.

To obtain positive critical measure, give an edge $v\to v'$ probability
$w_{v'}/(\lambda w_v)$. The path-cylinder mass telescopes to
$w_{v_m}/(\lambda^m w_{v_0})\le C\lambda^{-m}$. Fix one marked edge.
From every vertex there is a path crossing it in at most some fixed
number $L_0$ of steps, with probability at least $p>0$. Conditioning
successively in blocks of length $L_0$ bounds the probability of
avoiding the edge for $m$ blocks by $(1-p)^m$. Applying this after each
starting time shows that almost every path crosses a marked edge
infinitely often.

Feed the edge labels into the defect recurrence, with the two
initial integers zero and with starting index $N>8k$ sufficiently
large. The construction in the proof of Theorem~\ref{thm:gauge} gives
a point $x=\lim M_j/Q_j$ with
$|Q_jx-M_j|\le 4kQ_j/Q_{j+1}<1/8$ for $j\ge N$.
Thus the $M_j$ are its nearest integers and the prescribed labels
are its defects. Infinitely many marked labels give $\Lam(x)=k/4$
by the defect asymptotic. The pushforward of the path probability
therefore gives full mass to $S_k(G)$.

If a set of diameter $r$ satisfies $r_{n+1}\le r<r_n$, the difference
between two of its coded values of $M_j$ is less than
$1/4+Q_jr$. For $j<n$ with $j\ge N$ this is less than one, and
at $j=n$ there are at most three possible adjacent integers.
Consequently the set meets at most three label cylinders of length
$n-N$. From a prescribed start vertex each label word specifies at
most one path. Thus its mass is at most $3C\lambda^{-(n-N)}$, which is
bounded by a fixed multiple of $h_\lambda(r)$ for the interpolated gauge
$h_\lambda(r_n)=\lambda^{-n}$. Singleton mass is zero by the same
cylinder estimate. The construction is in a real interval of length
less than $1/2$ if $N$ is large, so projection to the circle is
isometric there. The mass-distribution bound gives
$\mathcal H^{h_\lambda}(S_k(G))>0$.

For Borel measurability, intersect with the Borel level $S_k$. For each
start and vertex, the condition that every finite prefix is admissible
is a countable intersection of Borel conditions on the rounding
functions. An infinite path exists by finite branching. Take the
countable union over starts and the finite union over vertices.
For $y\in L$, the defect asymptotic gives $d_j(y)=0$ eventually.
Because the nearest-integer errors of $x\in S_k$ and of $y$ tend to
zero, their nearest integers add eventually, so $d_j(x+y)=d_j(x)$
eventually. Hence $S_k(G)$ is invariant under the dense group $L$
and is Lebesgue-null. If its positive critical measure were finite,
its normalized restriction would be invariant under this dense group
and, by continuity against continuous test functions, under the
entire circle. It would then be Haar measure, a contradiction.
For a nonempty open set $U$, the family $U+\ell$, $\ell\in L\bmod 1$,
covers the circle by density and has a finite subcover by compactness.
Invariance under those translations would make the whole level have
finite measure if $S_k(G)\cap U$ had finite measure. This proves local
infinitude. The comparison of adjacent scales in
Corollary~\ref{cor:liminf-gauges} now gives the three conclusions for
every gauge. The middle case combines the countable covering by
fixed initial-pair classes with local infinitude on the full level.

For a general finite graph, each infinite path eventually remains
in one strongly connected component, since the component graph is
acyclic. If its point is in $S_k$, its defects have absolute value
$k$ infinitely often. Thus the terminal component is marked, and
\[
 S_k(G)=\bigcup_{C\ {\rm marked}} S_k(C).
\]
For a marked component with $\lambda_C=1$, each vertex has exactly
one internal outgoing edge. Indeed, two outgoing edges at one vertex
give two return loops with distinct first edges. Traversing them in
the two orders gives two distinct loops of equal length. Concatenating
those equal-length loops contradicts the bounded path count at
Perron root one. A marked component is therefore a directed cycle
in this case. Each fixed start index, vertex and integer pair
then specifies a unique defect tail and at most one real point:
the recurrence determines all later $M_j$, and $M_j/Q_j\to x$.
This gives countability.

Finally, when $1<\lambda_C<\lambda_*$ and
$\rho_{\lambda_*}(g)<\infty$, a bounded subsequence at base $\lambda_*$
has normalized cost tending to zero at base $\lambda_C$ by the
factor $(\lambda_C/\lambda_*)^n$. Such components are null.
Countable components are also null. At least one component attains
$\lambda_*$, giving local infinitude or local non-$\sigma$-finiteness
in the corresponding regimes. The finite union of the component
covers gives the asserted global finite-cover conclusion.
\end{proof}

\begin{remark}[Two graph examples]
Forbidding successive marked labels gives the adjacency matrix
\[
 \begin{pmatrix}2k-1&2\\2k-1&0\end{pmatrix}.
\]
For $k\ge2$ its characteristic polynomial has opposite signs at $2$
and $2k+1$, so its Perron root is strictly between those values.
For a reducible example, take a disjoint unmarked vertex with the
$2k-1$ unmarked loops and a marked vertex with the two loops labelled
$\pm k$. For $k\ge2$ the larger adjacency root comes from the unmarked
vertex, while $S_k(G)$ has critical base $2$. Thus the root of the
whole adjacency matrix need not determine the gauge of an exact level.
\end{remark}

For comparison, the graph-directed similarity results of Mauldin
and Williams \cite[Theorems~3 and~4]{MW} use a spectral radius of weighted
similarity matrices and study the full construction. Here the scales
are supplied by $Q_n$, and the insistence on infinitely many marked
edges arises from exact-level membership. The local infinitude uses
the arithmetic translation group of the level.

\paragraph{Lean status.}
All three regimes of Theorem~\ref{thm:finite-state-tail} are now kernel-checked at Layer~A in
\texttt{EpiFormal/\allowbreak Finite\allowbreak State\allowbreak Complete.lean} ($12$ audited declarations
ending in \texttt{thm\_\allowbreak finite\_\allowbreak state\_\allowbreak tail\_\allowbreak complete} on $\R$
and \texttt{thm\_\allowbreak finite\_\allowbreak state\_\allowbreak tail\_\allowbreak circle\_\allowbreak complete}
on \texttt{UnitAddCircle}, building on \texttt{Finite\allowbreak State\allowbreak Gauge.lean},
\texttt{Finite\allowbreak State\allowbreak Tail.lean}, and
\texttt{Finite\allowbreak State\allowbreak Positive.lean}, with base axioms only and zero \texttt{sorryAx}).
Specifically, under the deterministic outgoing-label hypothesis of this subsection, the formalization proves
compactness and subset embedding of the deterministic path tree (\texttt{graphTree}),
direct finite-sum Frostman positivity (\texttt{mkMetric\_\allowbreak graphTree\_\allowbreak pos}),
unmarked-tail nullity via backward-reachability stabilization at $L=|V|$ (\texttt{mkMetric\_\allowbreak unmarkedTail\_\allowbreak null}),
$L$-translation invariance of exact graph tails, and the arbitrary-gauge $\liminf$ comparison;
the positive right Perron--Frobenius eigenvector $(w_v,\lambda_C)$ on each marked strongly connected component is taken as an explicit hypothesis rather than derived via Lean matrix-eigenvalue machinery.

\subsection{General polynomial-growth denominator recurrences}\label{sec:poly-rec}

The quantization, exact gauge, three-way positive-level splitting, and
subcritical dimension phase transition extend from the Euler recurrence
\eqref{eq:rec} to arbitrary polynomial-growth three-term recurrences. Throughout
this subsection, fix a degree $r\in\Z_{\ge1}$, a leading coefficient
$c\in\Z_{\ge1}$, and an integer shift $d\in\Z$. Set $A_j:=cj^r+d$, choose an
integer $J_0\ge1$ such that $A_{J_0}=cJ_0^r+d\ge2$ (so $A_j$ is strictly
increasing and $A_j\ge2$ for all $j\ge J_0$), and let $(Q_j)_{j\ge0}$ be any
sequence of positive integers satisfying
\begin{equation}\label{eq:poly-rec}
Q_{j+1}=(cj^r+d)Q_j+Q_{j-1}\qquad\text{for all }j\ge J_0.
\end{equation}
For $x\in\R$ and $\alpha\ge0$, define the weighted approximation exponents
\[
\Lam^{(r,c,d)}(x):=\limsup_{j\to\infty}j^r\|Q_jx\|\in[0,\infty],\qquad
\Lam_\alpha^{(r,c,d)}(x):=\limsup_{j\to\infty}j^\alpha\|Q_jx\|\in[0,\infty],
\]
so that $\Lam^{(r,c,d)}=\Lam_r^{(r,c,d)}$. Fix the Borel rounding rule
$M_j(x):=\lfloor Q_jx+1/2\rfloor\in\Z$ and signed remainder
$u_j(x):=Q_jx-M_j(x)\in[-1/2,1/2)$ (so $|u_j(x)|=\|Q_jx\|$), and for $j\ge J_0$
define the integer recurrence defect (digit)
\begin{equation}\label{eq:poly-defect-id}
d_j(x):=M_{j+1}(x)-A_jM_j(x)-M_{j-1}(x)=A_ju_j(x)+u_{j-1}(x)-u_{j+1}(x)\in\Z.
\end{equation}

\begin{lemma}[Basic growth and non-asymptotic factorial bound, Layer~A in Lean]\label{lem:poly-growth}
Put $J_1:=\max(J_0+1,\,2|d|+1)$. Then:
\begin{enumerate}[label=(\roman*),leftmargin=2em]
\item For all $j\ge J_0$, $Q_{j+1}>A_jQ_j\ge2Q_j$.
\item For all $j\ge J_0+1$, $0<Q_{j-1}<\frac12Q_j<Q_j$, and hence
\[
cj^r+d<\frac{Q_{j+1}}{Q_j}<cj^r+d+1\le(c+|d|+1)j^r.
\]
\item For all $j\ge J_1$, $A_j\ge\frac{c}{2}j^r$, and for every $m\ge J_1$,
\[
Q_m\ge C_{J_1}\left(\frac{c}{2^{r+1}}\right)^m(m!)^r,\qquad
C_{J_1}:=Q_{J_1}\left(\frac{c}{2}\right)^{-J_1}\bigl((J_1-1)!\bigr)^{-r}>0.
\]
\item As $n\to\infty$, $\log Q_n=rn\log n+(\log c-r)n+O(\log n)$.
\end{enumerate}
\end{lemma}

\begin{proof}
Since $Q_{j-1}\ge1$ and $A_j\ge2$ for $j\ge J_0$, \eqref{eq:poly-rec} yields
$Q_{j+1}>A_jQ_j\ge2Q_j$, proving~(i). If $j\ge J_0+1$, then $j-1\ge J_0$, so~(i)
gives $Q_j>2Q_{j-1}$, i.e.\ $0<Q_{j-1}/Q_j<1/2<1$; adding $A_j$ gives
$cj^r+d<Q_{j+1}/Q_j<cj^r+d+1\le(c+|d|+1)j^r$ (using $j^r\ge1$), proving~(ii). For
$j\ge J_1$, we have $|d|<j/2\le\frac12 cj^r$ and therefore $A_j\ge\frac{c}{2}j^r$.
Multiplying $Q_{j+1}>A_jQ_j\ge\frac{c}{2}j^r Q_j$ over $j=J_1,\dots,m-1$ gives
\[
Q_m\ge Q_{J_1}\left(\frac{c}{2}\right)^{m-J_1}
\left(\frac{(m-1)!}{(J_1-1)!}\right)^r
=C_{J_1}\left(\frac{c}{2}\right)^m\frac{(m!)^r}{m^r}
\ge C_{J_1}\left(\frac{c}{2^{r+1}}\right)^m(m!)^r,
\]
since $m\le2^m$ for all $m\ge1$ implies $m^r\le(2^r)^m$, proving~(iii). Finally,
by~(ii), $\log(Q_{j+1}/Q_j)=r\log j+\log c+O(j^{-r})$ for $j\ge J_1$; summing
over $j=J_1,\dots,n-1$ and using $\sum_{j=1}^{n-1}\log j=n\log n-n+O(\log n)$
together with $\sum_{j=J_1}^{n-1}j^{-r}=O(\log n)$ (for $r=1$) or $O(1)$ (for
$r\ge2$) gives~(iv).
\end{proof}

When no pre-existing convergent sequence of $\ee$ is present, the role of the
group $L=\Z+\Z\frac{\ee+1}{2}$ is played by an intrinsic rank-$2$ lattice built
directly from the homogeneous recurrence $X_{j+1}=A_jX_j+X_{j-1}$ ($j\ge J_0$).
Let $(U_j)_{j\ge J_0-1}$ and $(V_j)_{j\ge J_0-1}$ be the two fundamental integer
solutions specified by
\[
U_{J_0-1}=1,\quad U_{J_0}=0;\qquad V_{J_0-1}=0,\quad V_{J_0}=1.
\]

\begin{lemma}[Intrinsic zero-level lattice and Casorati rigidity, Layer~A in Lean]\label{lem:poly-lattice}
\begin{enumerate}[label=(\roman*),leftmargin=2em]
\item For all $j\ge J_0-1$, $Q_j=Q_{J_0-1}U_j+Q_{J_0}V_j$ and
$U_jV_{j+1}-U_{j+1}V_j=(-1)^{j-J_0+1}$.
\item The limits $\alpha_U:=\lim_{j\to\infty}U_j/Q_j$ and
$\alpha_V:=\lim_{j\to\infty}V_j/Q_j$ exist, satisfy
$Q_{J_0-1}\alpha_U+Q_{J_0}\alpha_V=1$, and obey
\[
\left|\alpha_U-\frac{U_j}{Q_j}\right|<\frac{Q_{J_0}}{Q_jQ_{j+1}},\qquad
\left|\alpha_V-\frac{V_j}{Q_j}\right|<\frac{Q_{J_0-1}}{Q_jQ_{j+1}}
\qquad\text{for all }j\ge J_0.
\]
\item Set $g:=\gcd(Q_{J_0-1},Q_{J_0})\ge1$, write $Q_{J_0-1}=gq_0$ and
$Q_{J_0}=gq_1$ with $\gcd(q_0,q_1)=1$, pick integers $p_0,p_1\in\Z$ with
$p_0q_1-p_1q_0=1$, and put $\tilde Q_j:=Q_j/g=q_0U_j+q_1V_j$ and
$\tilde P_j:=p_0U_j+p_1V_j$. Then
\[
\theta:=\lim_{j\to\infty}\frac{\tilde P_j}{\tilde Q_j}=g(p_0\alpha_U+p_1\alpha_V)\in\R\setminus\Q,
\]
and the intrinsic rank-$2$ lattice
\[
L^{(r,c,d)}:=\Z\alpha_U+\Z\alpha_V=\frac{1}{g}\bigl(\Z+\Z\theta\bigr)
\]
is independent of the choice of $(p_0,p_1)$, contains $\Z$, and projects to a
dense countable subgroup $\bar L^{(r,c,d)}:=L^{(r,c,d)}/\Z$ of $\R/\Z$.
\item For $x\in\R$, $\Lam^{(r,c,d)}(x)=0$ if and only if $u_j(x)\to0$ and
eventually $d_j(x)=0$, if and only if $x\in L^{(r,c,d)}$. Moreover, every
$x\in L^{(r,c,d)}$ satisfies $\|Q_jx\|=O(1/Q_{j+1})$, whence
$\Lam_\alpha^{(r,c,d)}(x)=0$ for all $\alpha\ge0$.
\end{enumerate}
\end{lemma}

\begin{proof}
Part~(i) holds at $j=J_0-1$ by inspection and propagates to all $j\ge J_0-1$
because $(U_j),(V_j),(Q_j)$ satisfy the same recurrence \eqref{eq:poly-rec} whose
transfer matrices $\bigl(\begin{smallmatrix}A_j&1\\1&0\end{smallmatrix}\bigr)$
have determinant $-1$. Expanding $Q_j=Q_{J_0-1}U_j+Q_{J_0}V_j$ in the adjacent
ratio differences yields, for $j\ge J_0$,
\[
\frac{U_{j+1}}{Q_{j+1}}-\frac{U_j}{Q_j}
=\frac{(-1)^{j-J_0}Q_{J_0}}{Q_jQ_{j+1}},\qquad
\frac{V_{j+1}}{Q_{j+1}}-\frac{V_j}{Q_j}
=\frac{(-1)^{j-J_0+1}Q_{J_0-1}}{Q_jQ_{j+1}}.
\]
By Lemma~\ref{lem:poly-growth}(i), $Q_{j+2}>2Q_{j+1}>Q_j$ for all $j\ge J_0$,
so $1/(Q_jQ_{j+1})$ decreases strictly to $0$; the alternating series test
gives convergence and bounds the tail strictly by the magnitude of the first
term, proving~(ii) (and $Q_{J_0-1}\alpha_U+Q_{J_0}\alpha_V=1$ follows by taking
$j\to\infty$ in $Q_{J_0-1}(U_j/Q_j)+Q_{J_0}(V_j/Q_j)=1$).

For~(iii), $\bigl(\begin{smallmatrix}p_0&p_1\\q_0&q_1\end{smallmatrix}\bigr)\in\mathrm{SL}_2(\Z)$
gives $\tilde P_j\tilde Q_{j+1}-\tilde P_{j+1}\tilde Q_j=(-1)^{j-J_0+1}$, so
$\theta=\lim_{j\to\infty}\tilde P_j/\tilde Q_j=\lim_{j\to\infty}g(p_0U_j+p_1V_j)/Q_j=g(p_0\alpha_U+p_1\alpha_V)$
satisfies $0<|\theta-\tilde P_j/\tilde Q_j|<(\tilde Q_j\tilde Q_{j+1})^{-1}$ for
all $j\ge J_0$. If $\theta=p/q\in\Q$ ($q\ge1$), then
$0<|q\tilde P_j-p\tilde Q_j|<q/\tilde Q_{j+1}<1$ for large $j$, contradicting
$|q\tilde P_j-p\tilde Q_j|\in\Z_{\ge1}$; hence $\theta\notin\Q$. Since
$\bigl(\begin{smallmatrix}p_0&p_1\\q_0&q_1\end{smallmatrix}\bigr)\in\mathrm{SL}_2(\Z)$
is unimodular and $q_0\alpha_U+q_1\alpha_V=1/g$, we have
$\Z\alpha_U+\Z\alpha_V=\Z(q_0\alpha_U+q_1\alpha_V)+\Z(p_0\alpha_U+p_1\alpha_V)
=\Z\frac{1}{g}+\Z\frac{\theta}{g}=\frac{1}{g}(\Z+\Z\theta)$, which is thus
independent of $(p_0,p_1)$ and contains $1=g\cdot\frac{1}{g}$. Because
$\theta/g\notin\Q$, Kronecker's theorem shows that $\bar L^{(r,c,d)}$ is dense
in $\R/\Z$.

For~(iv), if $\Lam^{(r,c,d)}(x)=0$, then $j^r|u_j(x)|\to0$, so
$|d_j(x)|\le cj^r|u_j(x)|+|d||u_j(x)|+|u_{j-1}(x)|+|u_{j+1}(x)|\to0$; since
$d_j(x)\in\Z$, eventually $d_j(x)=0$. Conversely, if $u_j(x)\to0$ and $d_j(x)=0$
for all $j\ge N\ge J_0+1$, then by~(i)
$\det\bigl(\begin{smallmatrix}U_{N-1}&V_{N-1}\\U_N&V_N\end{smallmatrix}\bigr)=\pm1$,
so $(M_{N-1},M_N)=a(U_{N-1},U_N)+b(V_{N-1},V_N)$ for unique $a,b\in\Z$, and by
induction $M_j(x)=aU_j+bV_j$ for all $j\ge N-1$; dividing by $Q_j$ and using
$u_j(x)\to0$ yields $x=a\alpha_U+b\alpha_V\in L^{(r,c,d)}$. Finally, if
$x=a\alpha_U+b\alpha_V\in L^{(r,c,d)}$, then~(ii) gives
$\|Q_jx\|\le|Q_jx-(aU_j+bV_j)|\le(|a|Q_{J_0}+|b|Q_{J_0-1})/Q_{j+1}$ for
$j\ge J_0$, and Lemma~\ref{lem:poly-growth}(iii) yields $j^\alpha/Q_{j+1}\to0$
for every $\alpha\ge0$.
\end{proof}

\begin{theorem}[General $(cj^r+d)$ recurrences: complete spectrum discretization,
Layer~A in Lean]\label{thm:general-rec}
For any $r,c\in\Z_{\ge1}$, $d\in\Z$, and $(Q_j)_{j\ge0}$ satisfying
\eqref{eq:poly-rec}:
\begin{enumerate}[label=(\roman*),leftmargin=2em]
\item Whenever $u_j(x)\to0$ (in particular, whenever $\Lam^{(r,c,d)}(x)<\infty$),
the digits $d_j(x)$ satisfy
\[
\lim_{j\to\infty}\bigl(d_j(x)-cj^ru_j(x)\bigr)=0,\qquad
\lim_{j\to\infty}\bigl(|d_j(x)|-cj^r\|Q_jx\|\bigr)=0.
\]
\item The full spectrum of $\Lam^{(r,c,d)}$ is completely discrete with step $1/c$:
\[
\Lam^{(r,c,d)}(\R)=\frac{1}{c}\Z_{\ge0}\cup\{\infty\}
=\left\{0,\frac{1}{c},\frac{2}{c},\frac{3}{c},\dots\right\}\cup\{\infty\},
\]
so every open interval $\bigl(\frac{k}{c},\frac{k+1}{c}\bigr)$ ($k\in\Z_{\ge0}$)
is an empty spectral gap, and for every $t\in\bigl[\frac{k}{c},\frac{k+1}{c}\bigr)$
the sublevel set $E_t^{(r,c,d)}:=\{x\in\R:\Lam^{(r,c,d)}(x)\le t\}$ equals
$E_{k/c}^{(r,c,d)}=\{x\in\R:u_j(x)\to0\text{ and }\limsup_{j\to\infty}|d_j(x)|\le k\}$
(with $E_t^{(r,c,d)}=L^{(r,c,d)}$ for $t<1/c$).
\end{enumerate}
(Lean: for all $r\ge1$, proved at Layer~A in
\texttt{EpiFormal/\allowbreak TrichotomyAndPolyExtension.lean} via
\texttt{PolyRecData.\allowbreak poly\_\allowbreak signed\_\allowbreak defect\_\allowbreak approx},
\texttt{poly\_\allowbreak abs\_\allowbreak defect\_\allowbreak approx},
\texttt{poly\_\allowbreak spectral\_\allowbreak gap\_\allowbreak open\_\allowbreak interval},
\texttt{poly\_\allowbreak levelSet\_\allowbreak eq\_\allowbreak on\_\allowbreak gap},
\texttt{poly\_\allowbreak levelSet\_\allowbreak below\_\allowbreak eq\_\allowbreak lattice}, and
\texttt{poly\_\allowbreak range\_\allowbreak apprExponent}; the $r=1$
specialization also appears in
\texttt{EpiFormal/\allowbreak Full\allowbreak Spectrum\allowbreak Discrete.lean}.)
\end{theorem}

\begin{proof}
Rearranging \eqref{eq:poly-defect-id} gives
$d_j(x)-cj^ru_j(x)=du_j(x)+u_{j-1}(x)-u_{j+1}(x)\to0$ whenever $u_j(x)\to0$
(which holds whenever $\Lam^{(r,c,d)}(x)<\infty$), and the reverse triangle
inequality $\bigl||d_j(x)|-cj^r|u_j(x)|\bigr|\le|d_j(x)-cj^ru_j(x)|$ proves~(i).
Since $|d_j(x)|\in\Z_{\ge0}$, whenever $\Lam^{(r,c,d)}(x)<\infty$ we obtain
$c\Lam^{(r,c,d)}(x)=\limsup_{j\to\infty}|d_j(x)|\in\Z_{\ge0}$, excluding all
open intervals $(k/c,(k+1)/c)$ and forcing $E_t^{(r,c,d)}=E_{k/c}^{(r,c,d)}$ on
$[k/c,(k+1)/c)$. Conversely, $0\in L^{(r,c,d)}$ realizes $0$,
Theorem~\ref{thm:poly-phase-transition}(i) below realizes $\infty$ on a set of
full Lebesgue measure, and for each $k\ge1$ the constant-digit construction
$d_j\equiv k$ in Theorem~\ref{thm:poly-gauge-trichotomy} below produces a point
$x_k^+$ with $|u_j(x_k^+)|<4k/A_j=O(j^{-r})\to0$ and $d_j(x_k^+)=k$ for large
$j$, giving $\Lam^{(r,c,d)}(x_k^+)=k/c$ by~(i). The $t<1/c$ identity follows from
Lemma~\ref{lem:poly-lattice}(iv).
\end{proof}

For each integer $k\ge1$, define the continuous nondecreasing critical gauge
$h_k^{(r,c,d)}:[0,\infty)\to[0,\infty)$ by $h_k^{(r,c,d)}(0)=0$,
$h_k^{(r,c,d)}(1/Q_n)=(2k+1)^{-n}$ for all $n\ge J_0$, linear interpolation on
each $[1/Q_{n+1},1/Q_n]$ ($n\ge J_0$), and constant value $(2k+1)^{-J_0}$ for
$t\ge1/Q_{J_0}$ (well-defined because $Q_n$ is strictly increasing for
$n\ge J_0$ by Lemma~\ref{lem:poly-growth}(i)). Similarly define the universal
binary gauge $g_2^{(r,c,d)}(1/Q_n):=2^{-n}$ ($n\ge J_0$).

\begin{theorem}[Critical gauges and positive-level trichotomy for general $(cj^r+d)$, Layer~A in Lean]\label{thm:poly-gauge-trichotomy}
Fix an integer $k\ge1$. At the positive spectral level $k/c$:
\begin{enumerate}[label=(\roman*),leftmargin=2em]
\item \textbf{Limsup level (alphabet size $2k+1$, gauge base $2k+1$):}
$x\in E_{k/c}^{(r,c,d)}$ if and only if $u_j(x)\to0$ and eventually
$d_j(x)\in\{-k,\dots,k\}$. On $\R/\Z$, the exact level
$S_k^{(r,c,d)}:=\{x\in\R/\Z:\Lam^{(r,c,d)}(x)=k/c\}$ has infinite but
$\sigma$-finite $\mathcal H^{h_k^{(r,c,d)}}$ measure (with infinite measure in
every nonempty open arc), whereas $\{\Lam^{(r,c,d)}<k/c\}$ is
$\mathcal H^{h_k^{(r,c,d)}}$-null.
\item \textbf{Unsigned exact limit set (universal binary alphabet $\{\pm k\}$, gauge base $2$):}
$E_{\mathrm{lim}}^{(r,c,d)}(k/c):=\{x\in\R:\lim_{j\to\infty}j^r\|Q_jx\|=k/c\}$
consists precisely of those $x$ with $u_j(x)\to0$ and eventually
$d_j(x)\in\{-k,+k\}$ (two branches per step, independently of $k\ge1$). On
$\R/\Z$, $E_{\mathrm{lim}}^{(r,c,d)}(k/c)$ has infinite, $\sigma$-finite
$\mathcal H^{g_2^{(r,c,d)}}$ measure (infinite on every nonempty open arc),
while $\mathcal H^{h_k^{(r,c,d)}}(E_{\mathrm{lim}}^{(r,c,d)}(k/c))=0$. For
$t\notin\frac1c\Z_{\ge0}$, $E_{\mathrm{lim}}^{(r,c,d)}(t)=\emptyset$.
\item \textbf{Signed exact limit sets (single tail branch, coset collapse):}
$E_{\mathrm{lim}}^{(r,c,d),\pm}(k/c):=\{x\in\R:\lim_{j\to\infty}j^ru_j(x)=\pm k/c\}$
consists precisely of those $x$ with $u_j(x)\to0$ and eventually $d_j(x)=\pm k$.
With $x_k^+$ built from $d_j\equiv+k$ ($j\ge N$) and $x_k^-:=-x_k^+$, each signed
exact limit set collapses to a single countable coset of the intrinsic lattice
$L^{(r,c,d)}$ of Lemma~\ref{lem:poly-lattice}:
\[
E_{\mathrm{lim}}^{(r,c,d),+}(k/c)=x_k^+ + L^{(r,c,d)},\qquad
E_{\mathrm{lim}}^{(r,c,d),-}(k/c)=x_k^- + L^{(r,c,d)},
\]
and is therefore null for every non-atomic Hausdorff gauge.
\end{enumerate}
\end{theorem}

\begin{proof}
The digit characterizations follow directly from Theorem~\ref{thm:general-rec}.
For the upper bounds and $\sigma$-finiteness on $[0,1]$ (hence on $\R/\Z$ since
the quotient map is $1$-Lipschitz), partition $E_{k/c}^{(r,c,d)}\cap[0,1]$ into
countably many Borel classes $\mathcal C(N,m_{N-1},m_N)$ ($N\ge J_0+1$,
$m_{N-1},m_N\in\Z$) requiring $M_{N-1}=m_{N-1}$, $M_N=m_N$,
$|u_j|\le(k+1)/(cj^r)$, and $d_j\in\Sigma$ for all $j\ge N$, where
$\Sigma=\{-k,\dots,k\}$ ($|\Sigma|=2k+1$) for $E_{k/c}^{(r,c,d)}$,
$\Sigma=\{-(k-1),\dots,k-1\}$ ($|\Sigma|=2k-1$) for $\{\Lam^{(r,c,d)}<k/c\}$,
and $\Sigma=\{-k,+k\}$ ($|\Sigma|=2$) for $E_{\mathrm{lim}}^{(r,c,d)}(k/c)$.
At level $n>N$, the recurrence $M_{j+1}=A_jM_j+M_{j-1}+d_j$ yields at most
$|\Sigma|^{n-N}$ values of $M_n$, each confining $x$ to an interval of diameter
at most $1/Q_n$ (using $|u_n|\le1/2$). Since $n>J_0$, covering with these
intervals gives $\mathcal H_\delta^{h_k^{(r,c,d)}}$ cost at most
$(2k+1)^{-N}\bigl(|\Sigma|/(2k+1)\bigr)^{n-N}$ (bounded by $(2k+1)^{-N}<\infty$
when $|\Sigma|=2k+1$, and $\to0$ when $|\Sigma|<2k+1$), and when $\Sigma=\{-k,+k\}$
gives $\mathcal H_\delta^{g_2^{(r,c,d)}}$ cost at most $2^{n-N}2^{-n}=2^{-N}<\infty$.

For the lower bounds, choose an integer $N\ge J_0+1$ such that
$A_N=cN^r+d>32k$. Since $A_j\ge A_N>32k\ge32>6$ and $Q_N\ge1$, Lemma~\ref{lem:poly-growth}(i)
gives $Q_{N+1}>32kQ_N\ge32k$ and $Q_{j+1}>6Q_j$ for all $j\ge N$. Set
$M_{N-1}=M_N=0$ and, for $\Sigma\in\{\{-k,\dots,k\},\{-k,+k\}\}$ and
$\omega=(d_N,d_{N+1},\dots)\in\Sigma^{\{j\ge N\}}$, define $M_{j+1}=A_jM_j+M_{j-1}+d_j$
($j\ge N$). Then $W_j:=M_{j+1}Q_j-M_jQ_{j+1}$ satisfies $W_{N-1}=0$ and
$W_j=d_jQ_j-W_{j-1}$ ($j\ge N$). Since $j\ge N\ge J_0+1$ implies $Q_j>2Q_{j-1}$,
induction gives $|W_j|\le2kQ_j$ for all $j\ge N-1$, whence
$|M_{j+1}/Q_{j+1}-M_j/Q_j|\le2k/Q_{j+1}$. Using $Q_{m+1+s}\ge2^sQ_{m+1}$,
$x(\omega):=\lim_{j\to\infty}M_j/Q_j$ satisfies $|x(\omega)-M_m/Q_m|\le4k/Q_{m+1}$
for all $m\ge N$ (in particular $|x(\omega)|\le4k/Q_{N+1}$ since $M_N=0$).
Multiplying by $Q_m$ gives $|Q_mx(\omega)-M_m|\le4kQ_m/Q_{m+1}\le4k/A_m\le4k/A_N<1/8<1/2$
for all $m\ge N$, and at $m=N-1$ we also have $Q_{N-1}|x(\omega)|\le4kQ_{N-1}/Q_{N+1}<1/16<1/2$
(since $Q_{N+1}>32kQ_N>64kQ_{N-1}$). Hence $M_m(x(\omega))=M_m$ for all $m\ge N-1$,
no half-integer ties occur, and $d_j(x(\omega))=d_j$ for all $j\ge N$. Moreover,
$K_\Sigma:=x(\Sigma^{\{j\ge N\}})\subseteq[-4k/Q_{N+1},4k/Q_{N+1}]$, an interval
of length $<1/4<1/2$ embedding isometrically into $\R/\Z$. Pushing forward the
uniform Bernoulli product measure on $\Sigma^{\{j\ge N\}}$ yields $\mu$ on $K_\Sigma$.
For any interval $U$ of length $1/Q_{n+1}\le h<1/Q_n$ ($n>N$) and any $x,y\in K_\Sigma\cap U$,
we have $|M_j(x)-M_j(y)|<\frac14+\frac{Q_j}{Q_n}\le\frac14+\frac16<1$ for $N\le j<n$
(so $M_j(x)=M_j(y)$) and $|M_n(x)-M_n(y)|<\frac14+Q_nh<\frac54<2$ (at most $3$ values).
Thus $K_\Sigma\cap U$ meets at most $3$ cylinders of depth $n-N$, so
$\mu(U)\le3|\Sigma|^{-(n-N)}=3|\Sigma|^{N+1}|\Sigma|^{-(n+1)}\le3|\Sigma|^{N+1}g_{|\Sigma|}(h)$
(as $n+1>J_0$), where $g_{2k+1}:=h_k^{(r,c,d)}$ and $g_2:=g_2^{(r,c,d)}$. The
mass distribution principle yields $\mathcal H^{g_{|\Sigma|}}(K_\Sigma)\ge(3|\Sigma|^{N+1})^{-1}>0$.
Taking $\Sigma=\{-k,\dots,k\}$ and subtracting the $\mathcal H^{h_k^{(r,c,d)}}$-null
set $\{\Lam^{(r,c,d)}<k/c\}$ gives $\mathcal H^{h_k^{(r,c,d)}}(S_k^{(r,c,d)})>0$;
taking $\Sigma=\{-k,+k\}$ gives $K_\Sigma\subseteq E_{\mathrm{lim}}^{(r,c,d)}(k/c)$
and $\mathcal H^{g_2^{(r,c,d)}}(E_{\mathrm{lim}}^{(r,c,d)}(k/c))>0$.

By Lemma~\ref{lem:poly-lattice}(iv), $j^r\|Q_jy\|\to0$ for every $y\in L^{(r,c,d)}$,
so both $S_k^{(r,c,d)}$ and $E_{\mathrm{lim}}^{(r,c,d)}(k/c)$ are Borel,
Lebesgue-null (by Theorem~\ref{thm:poly-phase-transition}(i), or directly since
$\operatorname{Leb}_{\R/\Z}(\{\|Q_jx\|\le m/j^r\})=2m/j^r\to0$), and invariant
under translation by the dense subgroup $\bar L^{(r,c,d)}\subset\R/\Z$ of
Lemma~\ref{lem:poly-lattice}(iii). The exact Haar-uniqueness and compactness
covering arguments of Theorem~\ref{thm:gauge} and Corollary~\ref{cor:exact-gauge}(i)
then force both $\mathcal H^{h_k^{(r,c,d)}}(S_k^{(r,c,d)}\cap I)=\infty$ and
$\mathcal H^{g_2^{(r,c,d)}}(E_{\mathrm{lim}}^{(r,c,d)}(k/c)\cap I)=\infty$ on
every nonempty open arc $I\subseteq\R/\Z$.

Finally, for~(iii), take $x_k^+$ from the constant sequence $d_j\equiv+k$ ($j\ge N$).
Since $|u_j(x_k^+)|<1/8<1/2$ for all $j\ge N-1$, $x_k^-=-x_k^+$ satisfies
$M_j(x_k^-)=-M_j(x_k^+)$ and $u_j(x_k^-)=-u_j(x_k^+)$ for $j\ge N-1$, so
$d_j(x_k^-)=-k$ for $j\ge N$ and $x_k^\pm\in E_{\mathrm{lim}}^{(r,c,d),\pm}(k/c)$.
If $x,y\in E_{\mathrm{lim}}^{(r,c,d),\pm}(k/c)$, then
$j^r\|Q_j(x-y)\|\le j^r|u_j(x)-u_j(y)|\to0$, so $x-y\in L^{(r,c,d)}$ by
Lemma~\ref{lem:poly-lattice}(iv). Conversely, if $z\in L^{(r,c,d)}$, then
$u_j(z)=o(j^{-r})\to0$ while $u_j(x_k^\pm)=O(j^{-r})\to0$, so eventually
$|u_j(x_k^\pm)+u_j(z)|<1/2$, forcing $u_j(x_k^\pm+z)=u_j(x_k^\pm)+u_j(z)$ and
$j^ru_j(x_k^\pm+z)\to\pm k/c$. Thus $E_{\mathrm{lim}}^{(r,c,d),\pm}(k/c)=x_k^\pm+L^{(r,c,d)}$.
\end{proof}

\begin{corollary}[Exact critical gauge for $(cj^r+d)$, Layer~A in Lean]\label{cor:poly-exact-gauge}
Let $k\ge1$ and let $g:[0,\infty)\to[0,\infty)$ be any nondecreasing gauge with $g(0)=0$.
\begin{enumerate}[label=(\roman*),leftmargin=2em]
\item If $g(1/Q_n)/h_k^{(r,c,d)}(1/Q_n)\to0$ as $n\to\infty$, then
$\mathcal H^g(\{\Lam^{(r,c,d)}\le k/c\})=0$.
\item If $h_k^{(r,c,d)}(t)/g(t)\to0$ as $t\to0^+$, then $S_k^{(r,c,d)}$ is not
$\sigma$-finite for $\mathcal H^g$ (in particular $\mathcal H^g(S_k^{(r,c,d)})=\infty$).
\end{enumerate}
The exact same dichotomy holds for $E_{\mathrm{lim}}^{(r,c,d)}(k/c)$ with
$g_2^{(r,c,d)}$ in place of $h_k^{(r,c,d)}$.
\end{corollary}

\begin{proof}
This follows by the identical covering comparison and $\sigma$-finiteness
contradiction of Corollary~\ref{cor:exact-gauge}(ii)--(iii), using the upper
covers from Theorem~\ref{thm:poly-gauge-trichotomy}.
\end{proof}

\begin{remark}[Interpolation-free logarithmic scale and the $1/r$ factor]\label{rem:poly-gauge-scale}
For $1/Q_{n+1}\le t\le 1/Q_n$ ($n\ge J_0$), $-\log h_k^{(r,c,d)}(t)$ lies between
$n\log(2k+1)$ and $(n+1)\log(2k+1)$. By Lemma~\ref{lem:poly-growth}(iv),
$\log(1/t)=rn\log n+O(n)$, so $\log\log(1/t)=\log n+O(\log\log n)\sim\log n$,
yielding $\frac{\log(1/t)}{\log\log(1/t)}\sim rn$ and therefore
\[
-\log h_k^{(r,c,d)}(t)\sim\frac{\log(2k+1)}{r}\cdot\frac{\log(1/t)}{\log\log(1/t)},\qquad
-\log g_2^{(r,c,d)}(t)\sim\frac{\log2}{r}\cdot\frac{\log(1/t)}{\log\log(1/t)}
\qquad(t\downarrow0).
\]
Consequently, any gauge $g$ whose logarithmic rate against
$\frac{\log(1/t)}{\log\log(1/t)}$ exceeds $\frac{\log(2k+1)}{r}$ gives
$\mathcal H^g(\{\Lam^{(r,c,d)}\le k/c\})=0$, while any strictly smaller rate
yields non-$\sigma$-finite $\mathcal H^g$ measure on $S_k^{(r,c,d)}$.
\end{remark}

\begin{theorem}[Critical dimension zero, subcritical phase transition $1-\alpha/r$,
and supercritical collapse, Layer~A in Lean]\label{thm:poly-phase-transition}
For any $r,c\in\Z_{\ge1}$, $d\in\Z$, and $(Q_j)_{j\ge0}$ satisfying \eqref{eq:poly-rec}:
\begin{enumerate}[label=(\roman*),leftmargin=2em]
\item \textbf{Critical weight ($\alpha=r$):} For every finite $t\in\R$,
$\dimH\{x\in\R:\Lam^{(r,c,d)}(x)\le t\}=0$. Thus
$F^{(r,c,d)}:=\{x\in\R:\Lam^{(r,c,d)}(x)<\infty\}$ has Hausdorff dimension $0$
and Lebesgue measure $0$, while its complement
$F_\infty^{(r,c,d)}:=\{x\in\R:\Lam^{(r,c,d)}(x)=\infty\}$ has full Lebesgue
measure and $\dimH(F_\infty^{(r,c,d)}\cap[a,b])=1$ for every $a<b$.
\item \textbf{Subcritical weights ($0\le\alpha<r$):} For every $0\le\alpha<r$ and
every $t\ge0$ (with $0\le t\le1/2$ when $\alpha=0$),
\[
\dimH\bigl\{x\in[0,1]:j^\alpha\|Q_jx\|\longrightarrow t\bigr\}
=\dimH\bigl\{x\in\R:\Lam_\alpha^{(r,c,d)}(x)\le t\bigr\}
=1-\frac{\alpha}{r}.
\]
(For $\alpha=0$ and $t>1/2$, the prescribed-limit set is empty while
$\{\Lam_0^{(r,c,d)}\le t\}=\R$ has dimension $1=1-0/r$.)
\item \textbf{Supercritical weights ($\alpha>r$):} For every finite $t\ge0$,
$\{x\in\R:\Lam_\alpha^{(r,c,d)}(x)\le t\}=L^{(r,c,d)}$ (countable, Hausdorff
dimension $0$), while for $t<0$ the sublevel set is empty.
\end{enumerate}
\end{theorem}

\begin{proof}
\emph{Step 1: Upper bound $\dimH\{x\in\R:\Lam_\alpha^{(r,c,d)}(x)\le t\}\le1-\alpha/r$
for all $0\le\alpha\le r$.} Fix $0\le\alpha\le r$ and $t\in\R$, and put
$t_0:=\max(t,0)+1>0$. By $\Z$-periodicity and countable stability, it suffices to
show $\dimH C_{J,\alpha}\le1-\alpha/r$ for each $J\ge J_1$ (from
Lemma~\ref{lem:poly-growth}), where
$C_{J,\alpha}:=[0,1]\cap\bigcap_{j=J}^\infty\{x:\|Q_jx\|\le t_0j^{-\alpha}\}$.
Put $r_{j,\alpha}:=t_0/(j^\alpha Q_j)$ and
$I_{j,\alpha}(k):=[k/Q_j-r_{j,\alpha},k/Q_j+r_{j,\alpha}]$. Any $x\in C_{J,\alpha}$
satisfies $x\in I_{j,\alpha}(k_j(x))$ for $k_j(x)=\mathrm{round}(Q_jx)\in\{0,\dots,Q_j\}$,
with $|k_j(x)/Q_j-k_{j+1}(x)/Q_{j+1}|\le r_{j,\alpha}+r_{j+1,\alpha}$. Multiplying
by $2Q_{j+1}$ and using Lemma~\ref{lem:poly-growth}(ii) together with
$1/(j+1)^\alpha\le1\le j^{r-\alpha}$ ($0\le\alpha\le r$, $j\ge1$), each parent
$k$ yields child centers $k'$ in an interval of length
\[
2(r_{j,\alpha}+r_{j+1,\alpha})Q_{j+1}
=2t_0\left(\frac{Q_{j+1}}{j^\alpha Q_j}+\frac{1}{(j+1)^\alpha}\right)
\le 2t_0(c+|d|+2)j^{r-\alpha},
\]
so each parent has at most $Kj^{r-\alpha}$ children, where $K:=2t_0(c+|d|+2)+1$.
Hence the surviving center set $S_{n,\alpha}\subseteq\{0,\dots,Q_{J+n}\}$ at
level $m:=J+n\ge J_1$ satisfies
$|S_{n,\alpha}|\le(Q_J+1)K^n(m!)^{r-\alpha}$, while
Lemma~\ref{lem:poly-growth}(iii) and $m^\alpha\ge1$ give
$\operatorname{diam}I_{m,\alpha}(k)=2t_0/(m^\alpha Q_m)\le2t_0/(C_{J_1}(c/2^{r+1})^m(m!)^r)\to0$.
For any $s>1-\alpha/r$, set $\beta:=rs-(r-\alpha)=r(s-(1-\alpha/r))>0$ and
$K_2:=K(2^{r+1}/c)^s>0$. Using $m!\ge(m/\ee)^m$, the $s$-dimensional covering
cost satisfies
\begin{align*}
\sum_{k\in S_{n,\alpha}}\bigl(\operatorname{diam}I_{m,\alpha}(k)\bigr)^s
&\le(Q_J+1)K^{-J}\left(\frac{2t_0}{C_{J_1}}\right)^s K_2^m(m!)^{-\beta}\\
&\le(Q_J+1)K^{-J}\left(\frac{2t_0}{C_{J_1}}\right)^s
\left(\frac{\ee K_2^{1/\beta}}{m}\right)^{\beta m}\longrightarrow0
\end{align*}
as $m\to\infty$. Thus $\mathcal H^s(C_{J,\alpha})=0$ for every $s>1-\alpha/r$,
giving $\dimH\{x\in\R:\Lam_\alpha^{(r,c,d)}(x)\le t\}\le1-\alpha/r$. Taking
$\alpha=r$ proves~(i) (with $\operatorname{Leb}(F^{(r,c,d)})=0$ and
$\dimH(F_\infty^{(r,c,d)}\cap[a,b])=1$ exactly as in Corollary~\ref{cor:infinite-level}).

\emph{Step 2: Lower bound $\dimH\{x\in[0,1]:j^\alpha\|Q_jx\|\to t\}\ge1-\alpha/r$
for $0\le\alpha<r$.} Fix $0\le\alpha<r$ and $t\ge0$ (with $0\le t\le1/2$ if
$\alpha=0$). Put $w_j:=j^{-\alpha}/\log(j+2)$ and $r_j:=w_j/Q_j$. Since
$r-\alpha>0$ and $c\ge1$, Lemma~\ref{lem:poly-growth}(ii) gives
$2w_j(cj^r+d-1)-1\sim2c\,j^{r-\alpha}/\log(j+2)$; choose $N\ge\max(J_1,3)$ large
enough that for all $j\ge N$, $w_j<1/8$, $tj^{-\alpha}+w_j\le1/2$ (when
$\alpha>0$ or $\alpha=0,t<1/2$),
\[
2w_j(cj^r+d-1)-1\ge\frac{j^{r-\alpha}}{\log(j+2)}\ge b_j:=\left\lfloor\frac{j^{r-\alpha}}{\log(j+2)}\right\rfloor\ge2,
\]
and $Q_N\ge2$ (so $I_N(1)\subseteq(0,1)$). Defining the clamped shifts
$\theta_j\in[w_j,1/2-w_j]$ and disjoint intervals
$I_j(k):=[(k+\theta_j-w_j)/Q_j,(k+\theta_j+w_j)/Q_j]$ as in
Theorem~\ref{conj:weight}, every $x\in I_j(k)$ obeys
$|j^\alpha\|Q_jx\|-t|\le2/\log(j+2)\to0$. By Lemma~\ref{lem:poly-growth}(ii)
and $w_{j+1}\le w_j$, the nesting count inside each $I_j(k)$ is at least
\[
2(r_j-r_{j+1})Q_{j+1}-1=2\left(w_j\frac{Q_{j+1}}{Q_j}-w_{j+1}\right)-1
>2w_j(cj^r+d-1)-1\ge b_j\ge2.
\]
Retaining $b_j$ children per parent from $I_N(1)$ gives
$M_n:=\prod_{j=N}^{n-1}b_j$ intervals at level $n\ge N$ and a compact subset
$K\subseteq\{x\in[0,1]:j^\alpha\|Q_jx\|\to t\}$ carrying the uniform Bernoulli
probability measure $\mu$ ($\mu(I_n(k))=M_n^{-1}$). Since
$\log b_j=(r-\alpha)\log j-\log\log(j+2)+O(1)$, we have
$\log M_n=(r-\alpha)n\log n+o(n\log n)$, while Lemma~\ref{lem:poly-growth}(iv)
gives $\log Q_{n+1}=rn\log n+O(n)$. Thus for any $0<s<1-\alpha/r$,
\[
\log\left(\frac{3}{M_nQ_{n+1}^{-s}}\right)
=s\log Q_{n+1}-\log M_n+\log3
=-r\left(\left(1-\frac{\alpha}{r}\right)-s\right)n\log n+o(n\log n)\longrightarrow-\infty,
\]
so $3/M_n\le Q_{n+1}^{-s}$ for all $n\ge N_s$. Any set $U_i$ with
$Q_{n_i+1}^{-1}\le h_i:=\operatorname{diam}(U_i)<Q_{n_i}^{-1}$ ($n_i\ge N_s$)
satisfies $h_i+2r_{n_i}<(1+1/4)Q_{n_i}^{-1}<2Q_{n_i}^{-1}$ and thus meets at
most $3$ retained intervals, yielding $\mu(U_i)\le3/M_{n_i}\le Q_{n_i+1}^{-s}\le h_i^s$.
The mass distribution principle gives $\mathcal H^s(K)\ge1$, hence
$\dimH\{x\in[0,1]:j^\alpha\|Q_jx\|\to t\}\ge1-\alpha/r$. Combined with Step~1,
this proves~(ii).

\emph{Step 3: Supercritical collapse ($\alpha>r$).} If $t<0$, the sublevel set
is empty because $\Lam_\alpha^{(r,c,d)}\ge0$. Fix $t\ge0$. If
$\Lam_\alpha^{(r,c,d)}(x)\le t$, then $\|Q_jx\|=O(j^{-\alpha})=o(j^{-r})$, so
$\Lam^{(r,c,d)}(x)=0$ and $x\in L^{(r,c,d)}$ by Lemma~\ref{lem:poly-lattice}(iv);
conversely, every $x\in L^{(r,c,d)}$ satisfies $\Lam_\alpha^{(r,c,d)}(x)=0\le t$
by Lemma~\ref{lem:poly-lattice}(iv), proving~(iii).
\end{proof}

\begin{remark}[Comparison with Bugeaud--Liao--Rams]\label{rem:poly-BLR}
For $\alpha=0$ the dimension assertion in Theorem~\ref{thm:poly-phase-transition}(ii)
also follows from Corollary~8 of \cite{BLR} (stated there without proof), since
$Q_{j+1}/Q_j\to\infty$, by taking
constant centres $y_j=t$ exactly as for the convergent denominators of~$e$.
For $0<\alpha<r$, the right-hand side of the dimension formula in
\cite[Theorem~7]{BLR}, evaluated formally at $f_j(x)=Q_jx$ and $\delta_j=\kappa j^{-\alpha}$ ($\kappa>0$, $j\ge1$),
is $\liminf_n(\log Q_n+\sum_{1\le j<n}\log\delta_j)/(\log Q_n-\log\delta_n)=1-\alpha/r$,
the value in~(ii). That theorem does not apply, however: its condition~(1.6) requires
$\log f'_{j+1}/\log f'_j\to\infty$, whereas here $\log Q_{j+1}/\log Q_j\to1$, and the
proof above does not use that condition. (Hypotheses as stated in arXiv:1710.03596v1, where the proof of Theorem~7 is
given as a sketch.)
\end{remark}

\paragraph{Lean status for Section~\ref{sec:poly-rec}.}
The entire general-$(r,c,d)$ theory of Section~\ref{sec:poly-rec} for
\emph{all} degrees $r\ge1$, leading coefficients $c\ge1$, and shifts $d\in\Z$
is kernel-checked at Layer~A (base axioms only: \texttt{propext},
\texttt{Classical.choice}, \texttt{Quot.sound}; zero \texttt{sorryAx}) across
four modules totaling $68$ audited declarations under \texttt{PolyRecData}:
\texttt{EpiFormal/\allowbreak TrichotomyAndPolyExtension.lean} ($37$
declarations), \texttt{EpiFormal/\allowbreak PolyGaugeAndPhaseTransition.lean}
($22$ declarations), \texttt{EpiFormal/\allowbreak PolyWeightedUpper.lean} ($1$
declaration), and \texttt{EpiFormal/\allowbreak PolyWeightedLower.lean} ($8$
declarations):
\begin{itemize}[leftmargin=1.8em]
\item Lemma~\ref{lem:poly-growth}(i)--(iv) and decay of $m^\alpha/Q_m$
(\texttt{poly\_\allowbreak Q\_\allowbreak ge\_\allowbreak factorial},
\texttt{poly\_\allowbreak pow\_\allowbreak div\_\allowbreak Q\_\allowbreak tendsto\_\allowbreak zero},
\texttt{poly\_\allowbreak log\_\allowbreak Q\_\allowbreak asymptotic});
\item Lemma~\ref{lem:poly-lattice}(i)--(iv) including the Casorati identity,
fundamental solution error bounds, irrationality $\theta\notin\Q$, and
zero-level lattice equivalence
(\texttt{poly\_\allowbreak casorati},
\texttt{alphaU\_\allowbreak error\_\allowbreak bound},
\texttt{alphaV\_\allowbreak error\_\allowbreak bound},
\texttt{poly\_\allowbreak alpha\_\allowbreak linear\_\allowbreak relation},
\texttt{theta\_\allowbreak irrational},
\texttt{polyApprExponent\_\allowbreak eq\_\allowbreak zero\_\allowbreak tfae},
\texttt{polyApprExponent\_\allowbreak eq\_\allowbreak zero\_\allowbreak iff\_\allowbreak mem\_\allowbreak latticePoly});
\item Theorem~\ref{thm:general-rec}(i)--(ii), establishing the full discrete
spectrum $\Lam^{(r,c,d)}(\R)=\frac{1}{c}\Z_{\ge0}\cup\{\infty\}$ for all
$r\ge1$
(\texttt{poly\_\allowbreak signed\_\allowbreak defect\_\allowbreak approx},
\texttt{poly\_\allowbreak abs\_\allowbreak defect\_\allowbreak approx},
\texttt{polyApprExponent\_\allowbreak xPlusPoly},
\texttt{poly\_\allowbreak spectral\_\allowbreak gap\_\allowbreak open\_\allowbreak interval},
\texttt{poly\_\allowbreak levelSet\_\allowbreak eq\_\allowbreak on\_\allowbreak gap},
\texttt{poly\_\allowbreak levelSet\_\allowbreak below\_\allowbreak eq\_\allowbreak lattice},
\texttt{poly\_\allowbreak range\_\allowbreak apprExponent});
\item Theorem~\ref{thm:poly-gauge-trichotomy}(i)--(iii),
Corollary~\ref{cor:poly-exact-gauge}, and Remark~\ref{rem:poly-gauge-scale} on
both $\R$ and $\R/\Z$ (\texttt{UnitAddCircle}): critical-gauge and
binary-gauge nullity, $\sigma$-finiteness, open-interval/open-arc infinitude,
fast/slow gauge dichotomy, $\liminf$ four-way characterizations, logarithmic
gauge scale (\texttt{polyInterpGauge\_\allowbreak log\_\allowbreak scale\_\allowbreak tendsto},
\texttt{mkMetric\_\allowbreak polyCritGauge\_\allowbreak polyCircleExactLevel\_\allowbreak inter\_\allowbreak open\_\allowbreak eq\_\allowbreak top},
\texttt{mkMetric\_\allowbreak polyBinaryGauge\_\allowbreak polyCircleUnsignedExactLimit\_\allowbreak inter\_\allowbreak open\_\allowbreak eq\_\allowbreak top},
\texttt{polyExactLevel\_\allowbreak fast\_\allowbreak gauge\_\allowbreak null},
\texttt{polyExactLevel\_\allowbreak slow\_\allowbreak gauge\_\allowbreak non\_\allowbreak sigma\_\allowbreak finite},
\texttt{polyExactLevel\_\allowbreak liminf\_\allowbreak four\_\allowbreak way},
\texttt{polyUnsignedExactLimit\_\allowbreak liminf\_\allowbreak four\_\allowbreak way}),
and signed exact-limit coset collapse
$E_{\mathrm{lim}}^{(r,c,d),\pm}(k/c)=x_k^\pm+L^{(r,c,d)}$
(\texttt{poly\_\allowbreak trichotomy\_\allowbreak part\_\allowbreak iii\_\allowbreak complete});
\item Theorem~\ref{thm:poly-phase-transition}(i)--(iii): critical weight
$\alpha=r$ (\texttt{poly\_\allowbreak dimH\_\allowbreak levelSet\_\allowbreak finite},
\texttt{poly\_\allowbreak dimH\_\allowbreak finiteLevel\_\allowbreak zero},
\texttt{poly\_\allowbreak volume\_\allowbreak finiteLevel\_\allowbreak zero},
\texttt{poly\_\allowbreak ae\_\allowbreak apprExponent\_\allowbreak eq\_\allowbreak top},
\texttt{poly\_\allowbreak dimH\_\allowbreak infiniteLevel\_\allowbreak inter\_\allowbreak Icc\_\allowbreak eq\_\allowbreak one},
\texttt{poly\_\allowbreak dimH\_\allowbreak infiniteLevel\_\allowbreak eq\_\allowbreak one}),
subcritical $0\le\alpha<r$ exact Hausdorff dimension $1-\alpha/r$
(\texttt{Poly\allowbreak Weighted\allowbreak Upper.\allowbreak dimH\_\allowbreak poly\allowbreak Weighted\allowbreak Exponent\_\allowbreak le},
\texttt{Poly\allowbreak Weighted\allowbreak Lower.\allowbreak dimH\_\allowbreak poly\allowbreak Exact\allowbreak Limit\_\allowbreak eq},
\texttt{Poly\allowbreak Weighted\allowbreak Lower.\allowbreak dimH\_\allowbreak poly\allowbreak Weighted\allowbreak Sublevel\_\allowbreak eq},
\texttt{Poly\allowbreak Weighted\allowbreak Lower.\allowbreak poly\allowbreak Zero\allowbreak Weight\allowbreak Sublevel\_\allowbreak eq\_\allowbreak univ},
\texttt{Poly\allowbreak Weighted\allowbreak Lower.\allowbreak poly\allowbreak Exact\allowbreak Limit\_\allowbreak zero\_\allowbreak eq\_\allowbreak empty}),
and supercritical collapse $\alpha>r$
(\texttt{poly\_\allowbreak supercritical\_\allowbreak levelSet\_\allowbreak eq\_\allowbreak lattice},
\texttt{poly\_\allowbreak supercritical\_\allowbreak dimH\_\allowbreak zero},
\texttt{poly\_\allowbreak supercritical\_\allowbreak empty\_\allowbreak neg}).
\end{itemize}
These kernel-checked theorems are additionally corroborated numerically by the
$400/50$-digit \texttt{mpmath} checks in
\texttt{verify\_\allowbreak poly\_\allowbreak recurrence\_\allowbreak 2b.py}.

